\chapter{基于规则的知识获取}
\label{chap:rule-based-knowledge-acquisition}

第~\ref{chap:knowledge-acquisition-human}章讨论的依据来自人的逐篇判断，可靠但昂贵。是否存在一类可以重复执行、几乎不增加边际人力的判断依据？一个自然的想法是：把人判断时反复用到的条件直接写成程序规则，让机器自动执行。本章研究这种\emph{由规则提供依据}的知识获取方式，先从一个能暴露其全部难点的例子开始。

假设我们要从中文新闻中自动找出任职变动，并记录人员与组织之间的任职关系。最容易想到的规则是：句子里只要出现“任命”这个词，就判定它报告了任职变动。但拿三条新闻一测，这条规则会同时命中三种截然不同的情形：“海岬科技任命林舟为首席架构师”直接陈述了一次任职；“海岬科技未任命林舟”虽然出现了“任命”，却说的是\emph{没有}这回事；“关于任命林舟的传言已被否认”中，任命只是被转述的传言，而且随后被否认。三句话都含“任命”，却只有第一句真正支持目标事件。这个对照说明一个关键区分：程序在字符串里找到词语，只完成了\emph{字面检查}；要得到可用的标签，还必须再判断两件事——这条规则在当前语境下究竟\emph{应不应当触发}，以及一旦触发，推出的结论\emph{是否符合任务定义}。本章的大部分内容，都是在回答怎样把这两层判断也写进规则、并留下可供复核的依据。

要让上述判断稳定、可重复地执行，而不是每次都靠人临场决定，我们把判断所依据的东西显式写成\emph{规则}：有的规则规定“检查样本是否满足某些条件”，有的规定“去查询已有的实体、概念或关系事实”，还有的规定“最终答案必须满足哪些约束”。按这些事先写明的明确依据，由程序为样本产生任务标签、候选范围或关系标签，就是\term{规则标注}（rule-based annotation）。

要特别注意规则标注的产物地位：它产生的仍是\emph{候选}标签或约束，并不会因为出自程序就自动正确，因此每条结果都必须能说明“凭什么产生”，以便复核。理解规则时，最好把两个容易混淆的问题分开。一个是\emph{规则怎样形成}：规则文本由谁给出——可以是领域专家直接编写，可以是从大量样本的重复结构中归纳，也可以由模型先生成候选再由人确认。另一个是\emph{规则执行时依据什么产生标签}：是检查样本自身的模式，是关联外部已有的实体与事实，还是求解答案之间的约束。这两个维度相互独立：模型完全可以提出一条“去查询实体关系库事实”的规则（来源是模型、依据是外部事实），专家也可以亲手写一条“只比较数值阈值”的规则（来源是专家、依据是样本模式）。正因为来源和依据可以自由组合，二者必须分开记录，出错时才能定位到底是规则写错了，还是规则依赖的事实出了问题。

要把规则接进第~\ref{chap:knowledge-acquisition-overview}章那套含六个字段的样本标注，先要把规则执行的结果分得足够细，因为“没有得到标签”背后有三种完全不同的情况：句子里根本没有出现任命线索；否定语境过于复杂、程序无法可靠分析；以及这篇材料压根不是任务要处理的中文新闻。把三者都存成同一个空值，后续就无法决定是该补规则、弃权还是直接排除。

为此，一条规则在逻辑上拆成两个判断部分。其一是\emph{适用前提}，回答“这条规则是否适用于当前样本”，例如样本必须是中文新闻语体、事件时间在任务口径内，记为$A_v(x)$，其中$x$是样本、$v$是规则版本。其二是\emph{触发条件}，回答“在适用的前提下，样本是否命中了规则关心的模式”，记为$C_v(x)$；条件满足后如何拼出标签则由结果构造$O_v(x)$负责。

这里立刻遇到普通“是/否”布尔逻辑不够用的问题：除了条件成立、不成立，还常有\emph{暂时无法确定}，比如某个实体识别器没有输出，程序无从判断角色条件。因此本章用三值$\{\mathsf{T},\mathsf{F},\mathsf{U}\}$分别表示条件成立、不成立和无法确定。多个条件组合时，需要一种与三值配套的合取规则，称为\term{强Kleene合取}（strong Kleene conjunction）：只要有任一项为$\mathsf{F}$，合取就为$\mathsf{F}$（一个已确认不成立的条件足以否决）；只有所有项都为$\mathsf{T}$，合取才为$\mathsf{T}$；其余情形（没有明确假值、但至少一项未知）为$\mathsf{U}$。这样规定的好处是：一个已确认的假值能直接决定“不触发”，不会因为另一个条件恰好未知而被掩盖成“无法确定”。

令$D_v(x)=A_v(x)\land C_v(x)$表示适用前提与触发条件的合取。在一次条件检查\emph{成功执行完毕}之后，规则落入哪个分支由下式给出：
\begin{equation}
  b_v(x)=
  \begin{cases}
    \text{范围外}, & A_v(x)=\mathsf{F},\\
    \text{触发}, & D_v(x)=\mathsf{T},\\
    \text{未触发}, & A_v(x)\ne\mathsf{F}\ \text{且}\ D_v(x)=\mathsf{F},\\
    \text{无法确定}, & D_v(x)=\mathsf{U}.
  \end{cases}
  \label{eq:rule-based-knowledge-acquisition-interface}
\end{equation}
式中的$b_v(x)$只回答“这次检查走到了四个分支中的哪一个”，它本身并不代替那条完整的六字段样本标注，各分支还要翻译成相应的状态。

若适用前提$A_v(x)=\mathsf{F}$，说明对象根本不在规则的处理范围内，记为“范围外”：语义状态是“不适用”，来源动作为空——规则并没有尝试回答它。在适用前提未被否决的前提下，其余三种结果都作为\emph{具名触发结果}保存在取得条件字段$s_{\mathrm{cond}}$里，以便区分：$D_v(x)=\mathsf{T}$为“触发”，$D_v(x)=\mathsf{F}$为“未触发”，$D_v(x)=\mathsf{U}$为“无法确定”。对于那些“条件一满足就能直接拼出任务内容”的规则，三种状态对应不同动作：触发时来源动作为“作答”，给出$O_v(x)$，语义状态为已给出；未触发时规则只是说“我没看到目标模式”，因此没有来源动作、任务内容仍为未知（这一点对后文区分负类至关重要）；合取无法确定时，以“弃权”动作把内容保留为未知。

还要区分一类更复杂的规则：有些规则触发后并不立刻产出答案，而是进入一段后续求解（比如本章第三节要讲的约束推理）。对这类规则，“触发”只代表\emph{获准进入求解}，求解是否成功、最终来源动作是什么，必须另外记录，绝不能把“开始求解”误当成“已经作答”。

读者可能会把几种“没得到答案”的情形混为一谈，这里专门厘清执行状态。上述四个分支都建立在“条件检查这一\emph{动作本身}顺利跑完”的前提上，所以这一步的执行状态都是“成功”；但执行成功只表示程序正常工作，并不表示答案已被核验为真。如果后面还有求解步骤，求解有没有完成要单独保存其执行状态；而程序中途崩溃、网络中断属于技术故障，应记为“执行失败”，这不是规则在检查后主动选择弃权——二者绝不能混淆。

另一个贯穿全章的要点是负类。范围外、未触发、弃权这三种情况\emph{都没有证明样本属于负类}：规则没看到目标模式，不等于目标模式不存在。只有当一条规则明确检查负类定义、并据此返回一个负标签时，得到的样本标注才能当作负例。记住这一点，能避免把海量“规则没命中”的样本错误地当成反例喂给模型。

解决了“何时作答”，还要约定“规则返回的东西长什么样”，因为不同形态的结果需要的复核依据完全不同。规则结果$O_v(x)$至少有三种形态：它可以是一个\emph{候选标签}，比如“任职变动”；可以是一个\emph{候选范围}，只把答案缩小到$\{\text{管理岗位},\text{技术岗位}\}$这样的集合而不指定其一；还可以是一条\emph{关系约束}，比如“两个实体不得取同一身份”。

形态不同，复核时必须留下的证据也不同：靠样本自身匹配的规则，要留下命中的具体样本位置（哪个词、哪段字符）；靠查询外部事实的规则，要留下查到了哪条外部事实、来自哪个版本的库；靠推理的规则，则要留下完整的推理步骤。判断不出来时还要保存原因。表~\ref{tab:rule-based-knowledge-acquisition-record}沿用任命例子，把一条规则为了可复核而应记录的字段集中列出，读者可在后续三节中反复回到这张表。

\begin{table*}[htbp]
  \centering
  \begingroup
  \setlength{\tabcolsep}{4pt}
  \setlength{\parindent}{0pt}
  \setlength{\RaggedRightParindent}{0pt}
  \renewcommand{\arraystretch}{1.12}
  \begin{tabularx}{\linewidth}{>{\raggedright\arraybackslash}p{19mm}>{\raggedright\arraybackslash}p{32mm}>{\raggedright\arraybackslash}p{49mm}>{\raggedright\arraybackslash}X}
    \toprule
    字段 & 作用 & 教学取值 & 缺少时无法判断的问题 \\
    \midrule
    待检查对象 & 确定规则检查什么 & 一句中文新闻及其中的实体提及 & 规则究竟作用于整篇文档、句子还是实体对 \\
    适用前提 & 限定允许使用规则的范围 & 中文新闻语体，事件时间符合任务口径 & 对象在范围外，还是缺少判断范围所需的证据 \\
    触发条件 & 说明何时产生规则结果 & “任命”连接组织与人员，且不在否定或被否认的引语中 & 命中由哪个条件造成，新增例外改变了什么 \\
    规则结果 & 说明产生何种候选标签或约束 & 候选关系$(\text{林舟},\text{任职于},\text{海岬科技})$ & 返回的是类别、候选集合还是关系约束 \\
    无规则结果时的状态 & 区分范围、触发结果与动作 & 条件检查：范围外记不适用；未触发动作为\textemdash；合取无法确定时弃权 & 无任务答案是因为范围外、未触发，还是依据不足 \\
    命中证据 & 保留从待检查样本到规则结果的路径 & 关键词跨度、两个实体标识及局部依存路径 & 程序成功是否等于样本语义支持结论 \\
    适用范围及版本 & 支持更新与追溯 & 新闻语体；规则\texttt{appointment-v1}；实体关系库2026-09版 & 词表、解析器或实体属性或关系事实更新后哪些旧结果需要重算 \\
    \bottomrule
  \end{tabularx}
  \endgroup
  \caption{一条标注规则需要记录的内容。教学规则可以返回关系候选，也可以按任务改为类别或候选集合；字段的作用是把语义状态、来源动作、执行状态、规则依据和版本范围分开。}
  \label{tab:rule-based-knowledge-acquisition-record}
\end{table*}

沿任命例子继续想，规则为了得到结论实际上在求助三类不同的依据：它可能直接检查句子里的词语和角色（这是\emph{样本自身}就有的信息），可能去查明“林舟”对应现实中的哪个实体、岗位该归到哪个类别（这要借助\emph{样本之外已有的事实}），也可能不直接指定答案，而是规定几个候选答案之间必须满足什么关系（这是在\emph{答案空间}上做推理）。下面三节就按这三类依据展开：基于模式匹配的规则检查样本中的可观察条件，基于知识关联的规则使用已有的实体、概念与关系事实，基于约束推理的规则限制候选答案之间的合法组合。需要强调，一条真实规则常常同时使用这三类办法；但无论怎样组合，都必须为每一步单独留下依据，出错时才能判断问题究竟发生在匹配、关联还是求解环节。

% 标题、入口段与完整Snorkel边注共同检查剩余空间。
\Needspace{85mm}
\section{基于模式匹配的规则利用}
\label{sec:rule-based-knowledge-acquisition-pattern}

第一类依据完全来自样本自身。对一句待标注新闻，规则首先能利用的就是句中出现了哪些词语、它们在什么位置、数值是多少——这些都不需要求助样本之外的信息。规则可以检查这些可观察材料是否满足预先写好的条件，条件成立时给出候选结果；这种\emph{直接检查样本自身特征}的办法，称为\term{模式匹配}（pattern matching）。它是最基础、也最容易实现的一类规则依据。

模式匹配有一个被广泛使用的工程载体：Snorkel框架中的“标注函数”（labeling function）。一个标注函数就是一段可重复执行的程序，它依据字符串模式、启发式规则或外部资源，为送进来的样本给出一个类别标签，也可以在依据不足时明确放弃投票\scite{ratner2017}。把规则写成函数而非散落的脚本，好处是每条规则的触发条件、返回值和弃权行为都显式、可测试，也方便后续把成百上千条规则的结果统一融合（融合方法见第~\ref{chap:knowledge-fusion}章）。

这里要澄清一个术语差别，以免把Snorkel的约定和本章的状态字段错误对应。Snorkel原始工作把标注函数“不返回类别票”的所有情况统称为abstain（弃权），其中既包括“仔细看过但无法判断”，也包括“模式压根没匹配上”。但本章坚持把这两者分开：在条件检查阶段，\emph{明确未触发}（模式没命中、合取为假）时来源动作留空，只有\emph{无法确定而主动停止判断}时才记为弃权。第~\ref{chap:knowledge-acquisition-overview}章已经说明，“没看到目标”和“看到了但判断不了”是两件事，后续处理方式不同。

不过在送入标签融合模型的那一步，为了计算方便，可以把“未触发”和“主动弃权”这两种都没有类别票的结果，\emph{在融合层}映射成同一个“不投票”编码——前提是规则已成功执行、且对象确实处于共同适用范围内。必须强调这只是为融合产生的临时编码，它\emph{不覆盖}原始记录里的触发结果、来源动作、适用范围和弃权原因；这些字段仍要原样保留。而且，未请求、执行失败、适用范围外、或适用范围尚未对齐的结果，不属于“正常检查完但不投票”，必须另行处理，不能混进这个编码。如果将来还想分别为“未触发”和“主动弃权”建模、估计它们各自的发生规律，就需要在来源模型中保留条件状态并做相应扩展，此时融合层的单一“不投票”编码就不够用了。

在真正计算模式之前，通常还要做一道\emph{前处理}：把原始样本转换成可以直接检查的表示，例如把句子切分成词、把日志解析成字段值，等等；检查过程中命中的字符跨度、提取出的数值和各个中间条件的结果也都要保存下来。

之所以强调保存中间结果，是因为“规则没有给出可用答案”至少可能由三种不同障碍造成，必须能彼此区分。第一种是\emph{字段缺失}：前处理没能提供所需字段，条件根本无法求值，结果是无法确定。第二种是\emph{条件算出假值}：字段齐全、计算正常，但样本确实不满足条件，规则正常地未触发。第三种最隐蔽：\emph{条件成立却不支持目标语义}——字符串模式命中了，但命中处的真实含义并不是目标事件（正如“未任命”也含“任命”），这时产生的候选缺少语义依据，需要回头检查候选本身。前两种影响的是“触发/未触发/无法确定”的分支判断，第三种则发生在触发之后、针对已产生的候选；把它们一概记成“执行失败”会完全误导排查方向。下面先看符号条件和数值条件这两种最基本的条件，再看多个条件如何组合。它们全都依赖样本格式和前处理，一旦分词器或字段格式变化，原有规则必须重新验证。

\subsection{基于符号模式的规则利用}
\label{sec:rule-based-knowledge-acquisition-symbolic-pattern}

先看最直接的符号条件。一句新闻中的“任命”是一个清晰可观察的符号线索，但正如章首三例所示，只找到这个词还无法确定“谁任命了谁、任命是否真的发生”。要把判断做扎实，可以让符号条件由浅入深地分三层。

最浅的一层是\emph{关键词}：只检查某个词是否出现。它实现最简单，但表达力也最弱，无法约束词的位置和上下文。进一层是\emph{正则表达式}：它除了要求词出现，还能限制字符的相对位置和相邻形式，比如要求“任命”后面紧跟一个人名、前面是一个机构名。最深的一层是\emph{句式模板或局部结构模式}：它借助句法分析，检查若干成分（任命者、被任命者、谓词）是否以规定的语法关系同时出现，而不只是位置邻近。这三层用到的前处理越来越多，能区分的情形也越来越细。需要说明，这套思路并不限于文本：当检查对象从新闻换成日志或其他结构化观测时，被检查的“符号”可以变成某个字段的取值或某个事件类型标记，但判断依据仍然是\emph{当前样本内部的符号结构}，没有借助样本之外的事实——这正是模式匹配与下一节“知识关联”的分界。

把任命规则具体化。待检查样本限定为一个已经切分好的句子后，规则要识别句中的谓词及其参与者：谁是任命者（应当是一个组织）、谁是被任命者（应当是一个人员）。光有角色还不够，触发条件还要排除两种“假任命”：谓词不能受否定词支配（否则是“未任命”），任命也不能只出现在一句随后被报道者否认的引语里（否则是“传言被否认”）。在样本处于适用范围的前提下，只有所有条件都成立，规则才返回任职变动候选以及相应的关系三元组；只要有一个条件明确为假就不触发；没有明确假值、却仍有条件无法确定时才弃权。

以“海岬科技任命林舟为首席架构师”为例，一条可复核的规则不应只吐出一个“是”，而应同时记录：“任命”一词的字符跨度、“海岬科技”作为组织的跨度、“林舟”作为人员的跨度，以及三者在句法上的局部结构。有了这些记录，复核者才能回到原句逐一核对——这两个实体确实充当了规则声称的角色吗？谓词与它们的关系确实支持“任职”吗？如果规则只留下最终标签、不留命中位置，出错时就无从复查。

现在可以看清章首三例在多层条件下如何被区分。三条文本都能通过最浅层的关键词检查（都含“任命”），真正的差别在于任命是否被\emph{肯定地陈述}。加入“任命谓词不受否定词支配”这一条件后，直接否定句“海岬科技未任命林舟”被排除，因为否定词明确支配了谓词，该条件为假。但“关于任命林舟的传言已被否认”更麻烦：在引语\emph{内部}，词语可能仍是肯定形式，否定发生在引语\emph{之外}，由报道者对整段引述作出否认。因此规则不能只看句内否定词，还要识别引语的边界、并判断报道者是否否认了引语内容；当句法分析无法可靠确定这些时，对应条件应记为未知而非简单地判真或判假。

图~\ref{fig:rule-based-knowledge-acquisition-pattern-rule}把适用前提、谓词条件、角色条件和非否定条件的结果分别保存、再统一合取，读者应重点观察“未知项”和“假值”如何共同决定最终分支：一个明确假值足以使规则未触发，哪怕另有条件未知；只有在不存在假值时，未知才导致弃权。最后要破除一个想当然：增加条件确实会\emph{收窄}触发范围、减少明显误命中，但新条件本身（比如否定检测、引语识别）也可能判断错误，因此不能用“条件写得多”来论证“规则更准确”——条件数量与正确性之间没有等号。

\begin{figure*}[htbp]
  \centering
  \begin{tikzpicture}[
  x=1cm,
  y=1cm,
  font=\figurefont\footnotesize,
  panel/.style={draw=BookRule,line width=0.8pt,fill=white},
  text box/.style={draw=BookMuted,line width=0.8pt,fill=BookPaper,minimum height=10mm,text width=43mm,align=left,inner sep=2mm},
  condition/.style={draw=FigureModel,line width=1pt,fill=FigureModel!10,minimum height=10mm,text width=38mm,align=center,inner sep=1.5mm},
  output/.style={draw=FigureOutput,line width=1pt,fill=FigureOutput!7,minimum height=12mm,text width=39mm,align=center,inner sep=1.5mm},
  combine/.style={condition,draw=BookInk,fill=BookPaper},
  main/.style={knowledge arrow,draw=BookInk,line width=1.1pt,rounded corners=1mm}
]
  \fill[white] (0,0) rectangle (16.6,8);
  \draw[panel] (0,0) rectangle (5.15,8);
  \draw[panel] (5.4,0) rectangle (11.05,8);
  \draw[panel] (11.3,0) rectangle (16.6,8);

  \node[anchor=north west,font=\figurefont\bfseries] at (0.2,7.8) {(a) 相同字面命中};
  \node[text box,anchor=north west] at (0.25,7.1) {肯定：海岬科技\textbf{任命}林舟。\\\textcolor{BookTeal}{字面命中}};
  \node[text box,anchor=north west] at (0.25,4.95) {否定：海岬科技\textbf{未任命}林舟。\\\textcolor{BookTeal}{字面命中}};
  \node[text box,anchor=north west] at (0.25,2.8) {引用后否认：“\textbf{任命}林舟”的传言被否认。\\\textcolor{BookTeal}{字面命中}};
  \node[anchor=south,text=BookMuted,font=\figurefont\scriptsize,align=center] at (2.575,0.25) {匹配成功只说明找到了跨度};

  \node[anchor=north west,font=\figurefont\bfseries] at (5.6,7.8) {(b) 条件记录后统一合取};
  \node[condition] (scope) at (8.2,6.65) {$A_v$：适用前提\\记录$\mathsf{T}/\mathsf{F}/\mathsf{U}$};
  \node[condition] (keyword) at (8.2,5.35) {$K$：出现目标谓词\\记录$\mathsf{T}/\mathsf{F}/\mathsf{U}$};
  \node[condition] (roles) at (8.2,4.05) {$E$：组织与人员角色匹配\\记录$\mathsf{T}/\mathsf{F}/\mathsf{U}$};
  \node[condition] (context) at (8.2,2.75) {$N$：无否定或引语否认\\记录$\mathsf{T}/\mathsf{F}/\mathsf{U}$};
  \node[combine] (conjunction) at (8.2,1.0) {保留$A_v$并计算\\$D_v=A_v\land K\land E\land N$};
  \foreach \condition in {scope,keyword,roles,context} {
    \draw[main,-] (\condition.east) -- (10.7,0 |- \condition.east);
    \fill[BookInk] (10.7,0 |- \condition.east) circle (1pt);
  }
  \draw[main] (10.7,6.65) -- (10.7,1.95) -| (conjunction.north);

  \node[anchor=north west,font=\figurefont\bfseries] at (11.5,7.8) {(c) 分支与来源动作};
  \node[output] (triggered) at (14.1,6.65) {$D_v=\mathsf{T}$：触发、作答\\给出类别、范围或关系候选};
  \node[output,draw=BookMuted,fill=BookPaper] (untriggered) at (14.1,4.9) {$A_v\ne\mathsf{F},\ D_v=\mathsf{F}$\\未触发；动作：\textemdash\\任务内容未知};
  \node[output,draw=BookAmber,fill=BookAmber!8] (abstain) at (14.1,3.15) {$D_v=\mathsf{U}$：无法确定\\动作：弃权；任务内容未知};
  \node[output,draw=BookMuted,fill=white] (outside) at (14.1,1.4) {$A_v=\mathsf{F}$：范围外\\语义：不适用；动作：\textemdash};
  \coordinate (output-split) at (11.55,1.0);
  \draw[main,-] (conjunction.east) -- (output-split);
  \draw[main,-] (output-split) -- (11.55,6.65);
  \fill[BookInk] (output-split) circle (1pt);
  \foreach \result in {triggered,untriggered,abstain,outside} {
    \draw[main] (11.55,0 |- \result.west) -- (\result.west);
  }
\end{tikzpicture}

  \caption{从字面命中到带例外的组合规则。左面板中相同谓词出现在肯定、否定和被否认的引语里；中面板分别记录适用前提$A_v$、谓词条件$K$、角色条件$E$和非否定条件$N$，再计算$D_v=A_v\land K\land E\land N$，并列记录不要求并行执行。右面板按式~\eqref{eq:rule-based-knowledge-acquisition-interface}区分范围外、触发、未触发与无法确定；\textemdash{}表示来源动作为空。所有分支均以执行成功为前提。遇到未知仍须检查可独立求值的条件，例如$A_v=\mathsf{T}$且$(K,E,N)=(\mathsf{T},\mathsf{U},\mathsf{F})$时应未触发，而非弃权。}
  \label{fig:rule-based-knowledge-acquisition-pattern-rule}
\end{figure*}

前面逐层加码的条件，代价是对前处理的依赖也层层加深，这一点在维护时必须算清楚。字面关键词看似最简单，但它依赖文本规范化和分词是否正确，而且极易受多义词影响。带位置约束的正则表达式能规定局部顺序，却仍然难以处理否定成分与谓词分居两句的跨句否定。句式模板能检查实体角色和依存路径，可它又把可靠性押在了实体识别器和句法分析器的质量上。

由此得到一条规律：规则使用的语言结构越丰富，需要核验的中间结果就越多——关键词命中、分词边界、实体类型、依存弧，每一环都可能出错。这意味着维护一条复杂规则，绝不只是偶尔改改规则文本：词形变化、新的拼写方式、分词器升级、领域特有的句式、乃至解析器版本更新，都可能改变中间结果，从而让原本正确的规则静默失效。这些前处理的版本应像规则本身一样被记录和回归测试。

这些符号条件\emph{从哪里来}，是与“怎么执行”相连但不同的另一个问题，呼应章首对“规则来源”和“规则依据”的区分。专家可以凭借经验直接写出条件；也可以从语料中反复出现的结构归纳出模板；还可以让模型先生成一批候选条件。无论候选来自哪里，都不能直接上线，而要回到原始样本上同时检验支持情形和反例：它在哪些句子上正确触发、在哪些句子上误触发？误触发究竟是因为条件写得太宽，还是适用前提划得不对？

排查时尤其要利用保存下来的触发位置和解析结果，因为这能区分两类性质不同的错误：一类是\emph{规则解释错了}——前处理给出的结构是对的，但规则把这个结构错误地对应到了目标语义；另一类是\emph{前处理就错了}——句法分析或实体识别根本没给出正确结构，规则在错误的结构上判断，自然不会对。前者要改规则，后者要改或更换前处理，混为一谈就会修错地方。

最后要承认模式匹配的两类固有遗漏，即使规则本身执行得完全正确也无法避免。第一类是\emph{语义/时效上的过度触发}：规则准确匹配到了一则任命公告，但任务只收录“已经生效”的任职关系，而这则公告可能是尚未生效的任命预告；模式命中了，目标事实却还不成立，此时规则结果不符合任务定义，不能直接当成事实。第二类方向相反，是\emph{覆盖不足}：某句确实属于目标类，但因为句子写法特殊、解析依据不足，规则无法确定而弃权，于是这个本该作答的样本被漏掉了。

这两类问题提醒我们：规则执行正确（程序没报错、条件求值无误）与规则结果在任务意义上正确，是两个层面的事。前者属于本章的工程正确性，后者要靠第~\ref{chap:knowledge-quality-assessment-correction}章的统一质量检查来评估和修正，不能在写规则阶段就默认两者一致。

\subsection{基于数值条件的规则利用}
\label{sec:rule-based-knowledge-acquisition-numeric-pattern}

符号模式处理的是“出现什么词、结构如何”，但很多判断对象根本不是词语，而是一个有大小、可比较的量。要判断设备是否过热，检查新闻里有没有“烫”“热”这些字是靠不住的，真正该看的是传感器的温度读数。对于测量值、计数、比例或距离，规则先取得可比较的数值，再判断它是否落在指定阈值或区间内；这种办法称为\term{数值条件规则}（numeric-condition rule）。

必须从一开始就划清一条界限：数值比较是两个层面的事。“比较计算本身有没有算错”是一回事——82 是否大于 80，可以机械地、确定地验证；“这个阈值 80 对当前设备、当前任务是否合理”则是另一回事，它依赖物理特性和任务口径，无法靠比较运算自证。计算正确只能说明规则\emph{忠实于}给定的观测值和公式，完全不说明阈值选得对。这一节后半会专门处理阈值从哪来、何时需要怀疑它。

用一个教学规则看清数值判断里的各种边界。假设任务规定：设备读数不低于$80\,^\circ\mathrm{C}$时返回“高温候选”；落在$[78,80)\,^\circ\mathrm{C}$这个紧贴阈值的区间内时，因为测量误差可能让真值跨过临界线，规则选择弃权；低于$78\,^\circ\mathrm{C}$时这条规则不触发。每个观测值还应附带单位、采集时间和传感器状态，否则数值的含义不完整。

在其余前提都满足时：读数$82\,^\circ\mathrm{C}$明确触发；$79\,^\circ\mathrm{C}$落在缓冲区里，规则弃权；若连单位都缺失、无法完成有意义的比较，同样弃权而不是乱猜。单位换算是一个特别容易出错、又特别能说明“忠实观测值”重要性的环节：$176\,^\circ\mathrm{F}$换算成摄氏恰好是$80\,^\circ\mathrm{C}$，达到了\emph{包含端点}的下界$80$，所以应当触发。但如果程序偷懒、把数字$176$直接当成$176\,^\circ\mathrm{C}$去比较，比较运算照样会顺利执行（176当然大于80）、甚至也触发，可这个候选依据的是一个被错误理解的观测值——程序没报错，结论却不可信。这类“能跑但语义错”的故障，正是数值规则必须保存单位和换算过程的原因。

这个温度例子说明，一条数值规则远不止“写一个大于号”那么简单，它实际是一条处理链：先从原始观测\emph{计算出}可比较的特征，再\emph{统一单位或基准}，然后根据区间端点是包含还是排除，执行严格（$>$）或非严格（$\geq$）比较，最后把比较结果翻译成触发、未触发或无法确定。

当规则使用多个区间（如分档报警）时，还要明确两件事：每个端点归哪个区间（$80$ 到底算上一档还是下一档），以及区间之间是否允许重叠；重叠或留下空隙都会让某些数值对应多个或零个结果。对缺失值、传感器溢出（如读数恒为最大值）、以及根本不可比较的值，也要预先规定处理方式，不能让它们静默地落入某一档。例子中$[78,80)$这个“不确定缓冲区”不是随便画的：它应当来自测量误差范围、数值舍入范围或其他有证据支持的限制。如果误差来源未知，规则可以保守地弃权，但不能为了显得严谨，就把任意划定的一条缓冲带美其名曰“经过验证的置信区间”——没有误差依据的缓冲区只是主观裕量。

阈值的\emph{来源}决定了它的身份，也决定了当情况变化时该怎么对待它，这里有三种典型来源，不能混用。其一，阈值由任务定义直接规定，比如教学任务硬性规定$80\,\mathrm{^\circ C}$为高温线。这时阈值其实在\emph{定义标签口径}：什么叫“高温”完全由它说了算；一旦口径改变，按旧口径标注的样本就需要重新标注。其二，阈值来自某个领域标准（如行业安全规范）。这时要保存它适用的环境条件和标准版本，因为标准会修订、也只在规定环境下成立。其三，阈值是用少量样本“调”出来的经验值。这时必须保存用来调参的验证样本和整个调整过程，并且用\emph{没有参与调参}的新样本做独立检验，否则很容易把“在开发样本上恰好吻合”误当成“对新样本也成立”。

无论哪种来源，阈值都不是一劳永逸的。当监测对象的群体构成、采集设备型号或时间分布发生变化时，原本合适的阈值可能整体偏移，需要重新评估它对\emph{当前}对象是否仍然适用。

即使阈值数值已经确定，要拿到一个候选仍需完成一系列前置步骤：计算特征、查询参照基准、换算单位、处理缺失值，每一步都是潜在的错误来源。当不同地区、设备或人群使用不同基准时（比如不同的温度标准、不同的正常值范围），比较之前还必须先选对\emph{与观测设备和单位相匹配}的那个基准；拿 A 设备的读数去套 B 人群的基准，结论没有意义。

正因为错误可能来自链条上的不同环节——基准选错、特征算错、或把范围外的样本误交给规则——规则应在每一步都记录观测值、计算结果和所用基准版本。这样维护者面对一个错误候选时，才能沿链条定位是哪一环出了问题，而不是笼统地说“规则不准”。最后要把一项\emph{持续成本}计入：传感器会漂移、观测分布会随时间变化，因此对输入分布的监测不是一次性工作，而是规则上线后长期要做的维护。

\subsection{基于组合条件的规则利用}
\label{sec:rule-based-knowledge-acquisition-composite-pattern}

实际规则很少只检查单个条件，任命规则就是典型：它必须\emph{同时}确认谓词出现、实体角色正确、语境没有否定，才有充分依据返回候选。把这三项分别记为$K$（出现任命谓词）、$E$（实体角色匹配）、$N$（没有否定或被否认的引语），触发条件就是它们的合取
\[
  C_v(x)=K\land E\land N.
\]
各项仍然取章首规定的三值$\{\mathsf{T},\mathsf{F},\mathsf{U}\}$，并按强Kleene合取计算。这里值得重申三值在组合时的现实意义：如果实体识别器对某个片段根本没有预测输出，角色条件$E$是无法确定（$\mathsf{U}$），而\emph{不是}“实体类型不符合”（$\mathsf{F}$）——把“没看出来”误记成“看出来不是”，会错误地让规则从不触发，漏掉本应弃权待查的样本。

当任务里既有特殊情形又有一般情形时，多条规则之间还需要一套\emph{控制策略}：先试哪条？某条给出结果后还要不要继续看别的？这与单条规则内部的合取是两个层面的问题。

第一种组织是\term{规则列表}（rule list）：规则按优先级从上到下依次检查。本节采用一种保守策略——只有当某条规则被确定地判为“未触发”或“范围外”时，才继续往下试；一旦触发就作答并停止；一旦无法确定就弃权并停止。这种“遇未知即停”的保守性体现在：假设排在前面的特殊规则给出$D_v(x)=\mathsf{U}$，那么即使后面的一般规则本来能够触发，列表也\emph{不}采用一般规则的答案，因为特殊情形尚未被排除，贸然采用一般答案可能正好踩中特例。列表末尾若有一个“默认分支”想返回负类，条件必须非常严格：只有在列表的适用前提已确认、前面各项都\emph{确定地}未触发或范围外、且任务明定“这些都不命中即构成负类”时，默认分支才能给负标签；否则只能保留为未触发、动作为空。

第二种组织是\term{规则树}（rule tree）：沿不同条件分支向下走。本节同样规定，遇到未知时沿一个专门的弃权出口离开，而不是把未知默认成“否”。需要强调，列表和树在未知时“弃权并停止/走弃权出口”只是\emph{本节选择的控制策略}，并不是强Kleene合取本身强制规定的；换一种更激进的策略，未知处理可以不同。第三种组织则是互不设优先级的\term{规则集合}（rule set）：各条规则独立执行、互不短路，系统保留每一条的原始结果（包括各自的未知），把冲突与融合留给后面的多来源融合阶段。

逐项检查并不意味着每次都要把所有昂贵计算跑完。合取有一个可以利用的性质：一旦某一项已经确定为假，整个合取的结果就已被锁定为假，其余项不再影响结果，因此可以\emph{短路}——直接跳过它们。短路既有计算意义，也有正确性要求。

举例来说，若适用前提已确认成立、但最廉价的关键词检查$K$已经为假，那就不必再运行昂贵的句法解析去求$E$和$N$了，因为合取必为假，省下计算。但短路不能用错方向：当$K$为真、$E$未知时，规则\emph{还不能}弃权，因为$N$仍可独立求值，应当继续检查否定条件。具体地，若$(K,E,N)=(\mathsf{T},\mathsf{U},\mathsf{F})$，由于$N=\mathsf{F}$这个明确假值已经存在，合取为假，标注状态应是“未触发、动作为空”，而\emph{不是}弃权。请记住这个次序：\textbf{假值优先于未知}——只有在“不存在任何假值”的前提下，未知项才会让最终结果变成无法确定。短路时还要记录哪些条件是真正求过值、哪些只是被跳过，以便复核。

同一条“假值优先”的规则也适用于适用前提与触发条件之间，要注意两种情形的处理次序不同。如果适用前提未知、但触发条件已经确定为假，那么$D_v(x)$仍为假（前提未知不足以挽救一个已被否决的触发条件），可以确认“未触发”；但如果适用前提\emph{本身}已为假，则应优先按“范围外”处理，而不是“未触发”——前者说对象不归这条规则管，后者说对象归你管但没命中，语义不同。标注状态里应能区分一个条件究竟是：已经求出真值、仍然未知、还是因短路被跳过。

再看否定例外与正向模式冲突的情形：一旦否定条件$N$被确认不成立（即确实存在否定或否认），那么即使一般的词语模式命中了“任命”，规则也应记为未触发，同时把\emph{两处证据}都保留下来——既要保留“任命”的命中位置，也要保留否定依据，因为复核者需要看到这两者如何相互否决。

搭好一个多条件规则后，不能只看“它在开发样本上是否全对”就收工，还要逐项理解每个条件\emph{实际起了什么作用}。有两个简单有效的诊断动作。一是\emph{逐项删除}：每次去掉一个条件，看哪些新样本会因此进入触发范围，从而知道这个条件究竟在拦截什么。二是检查\emph{条件共现}：若两个条件几乎总是同时成立或同时不成立，其中一个可能是冗余的，白白增加维护成本和前处理依赖。

还要避免一个推理错误：给合取规则增加条件，通常确实会\emph{缩小}触发覆盖，但覆盖变窄绝不等于可靠性提高——新增条件自身也可能判断错误，它在挡住一部分误触发的同时，可能引入新的漏触发。对从样本归纳或由模型生成的条件组合，这些诊断同样必不可少；而且条件组合得越深、前处理链条越长，需要调试的路径和依赖就越多，复杂度本身就是一种成本。

最后把三个层次的“决策”理清，避免混用：规则列表/规则树的控制策略决定\emph{多条规则或多个分支之间}如何选择结果；合取短路决定\emph{同一条规则内部}哪些计算可以安全跳过；而规则集合中相互独立的规则则是另一回事——它们应当各自保留结果、不做选择。

这一区分在多规则冲突时尤其重要。如果多条独立规则对同一样本给出不同标签，而实现时简单地“以后执行的覆盖先执行的”或随手丢弃某些结果，冲突就被悄悄隐藏了，事后无从判断哪个标签其实更有依据。正确做法是：在规则层老老实实保留每一条规则的版本、证据和原始结果（包括分歧），不提前裁决；如何把这些可能冲突的标签综合成最终结论，是第~\ref{chap:knowledge-fusion}章多来源融合要解决的问题，规则层不应越俎代庖。

\section{基于知识关联的规则利用}
\label{sec:rule-based-knowledge-acquisition-association}

模式匹配的全部依据都在样本内部，但有些判断单凭样本里的文字怎么都做不出来。在任命新闻里，就算程序准确地圈出了“林舟”这两个字，也无法知道他指的是现实中哪一位同名的人——叫这个名字的作者可能有好几位；就算确认了头衔是“首席架构师”，也还要知道任务采用的岗位体系把这个头衔归到哪一类。这里缺的不是“再写一个更精巧的字面模式”，因为无论模式怎么写，样本里根本没有区分两位“林舟”的信息。缺的是\emph{样本与外部已有知识之间的一座桥}。

这就引出第二类依据：先确定样本中的提及对应到知识库中的哪个已有实体、概念或事实，再利用这层对应去产生候选标签或约束，称为\term{外部事实关联}（external fact association）。它与模式匹配的根本区别在于：模式匹配回答“样本自身满足什么可观察条件”，而外部事实关联回答“样本里的东西对应到了外部世界的哪个对象，以及那个对象有哪些已知属性”。

使用外部事实时，初学者容易以为“查到了就可以直接用”，其实查到之后还有两个必须分开回答的问题。第一个问题是这项事实\emph{本身是否成立、是否仍然有效}：实体的属性可能已经过时（人已离职、机构已更名），知识库中的关系可能有时间有效期。第二个问题是这项事实\emph{能否用于当前这个样本}：即便事实在知识库中为真，它也未必在当前这句话里被表达——两个人真实存在雇佣关系，不等于每一句同时提到他们的话都在谈论雇佣。把“事实为真”错误地等同于“本句表达了该事实”，正是后面远程监督要警惕的问题。

外部事实常组织成\term{本体}（ontology）：一个明确记录概念及其相互关系（如上下位、从属）的体系。查询本体时必须核对它采用的分类口径——同一个岗位词，在不同岗位体系里可能落到不同位置。表~\ref{tab:rule-based-knowledge-acquisition-association}按\emph{实体、概念、关系}三种查询对象，分别比较需要建立什么对应、以及错误容易发生在哪一步。需要说明，一条规则可以同时使用多种关联，但它们之间没有固定的先后顺序，不必机械地“先实体、再概念、后关系”，应按样本和任务的实际需要组织。

\begin{table*}[htbp]
  \centering
  \begingroup
  \setlength{\tabcolsep}{2.5pt}
  \setlength{\parindent}{0pt}
  \setlength{\RaggedRightParindent}{0pt}
  \renewcommand{\arraystretch}{1.12}
  \begin{tabularx}{\linewidth}{>{\raggedright\arraybackslash}p{17mm}>{\raggedright\arraybackslash}p{25mm}>{\raggedright\arraybackslash}p{26mm}>{\raggedright\arraybackslash}p{30mm}>{\raggedright\arraybackslash}p{29mm}>{\raggedright\arraybackslash}X}
    \toprule
    关联类型 & 样本线索 & 查询对象 & 对应条件 & 候选标签或约束 & 主要歧义 \\
    \midrule
    实体匹配 & 名称、别名、标识及上下文 & 实体标识及属性 & 提及指向该实体；属性适合当前时间与任务 & 类别、属性或查询入口 & 同名未必同实体；无候选与多候选有别 \\
    概念关联 & 词语、短语及局部用法 & 词典、本体或分类概念 & 表达指向该概念；目标属性允许沿层级继承 & 概念标签、上位类或候选范围 & 相似未必同义；上位类不能确定唯一细类 \\
    关系匹配 & 已消歧实体对及证据范围 & 有方向和时间的关系事实 & 实体与关系定义对齐；句级需表达证据，组级需另立假设 & 句级关系候选或提及组监督 & 共现未必表达；缺失未必为假 \\
    \bottomrule
  \end{tabularx}
  \endgroup
  \caption{三类外部事实关联的待检查样本、依据与错误位置。表中每一行都要求同时保存样本侧证据和实体、概念与关系库版本，避免只留下最终标签而无法判断错误发生在哪一步。}
  \label{tab:rule-based-knowledge-acquisition-association}
\end{table*}

\subsection{基于实体匹配的规则利用}
\label{sec:rule-based-knowledge-acquisition-entity-matching}

要查询新闻中“林舟”的属性，前提是先确定这个名字到底指谁——这是一切后续查询的入口。本章用\term{实体匹配}（entity matching）泛指把样本中的对象与已有实体对应起来的过程；在文本场景下，它的具体操作称为\term{实体链接}（entity linking），即把文中的一个提及（mention）对应到知识库里那个有唯一标识的实体对象。

实体链接带来一个关键转变：匹配成功后，规则处理的对象从一段有歧义的文字跨度，变成了一个可以直接查询的实体标识。因此它绝不能退化成单纯的字符串比较——“林舟”这两个字相同，不代表指向同一人。实体链接可利用的样本线索包括提及本身的文字跨度、名称、别名、显式标识符（如工号、网址）以及周围上下文；它的结果有四种可能，且含义各不相同：唯一确定一个实体、得到多个暂时无法区分的候选、一个候选都没有、或证据不足而弃权。

用一个同名的教学例子说明消歧如何进行。设实体库里有两位都叫“林舟”的作者：\texttt{E-017}的属性与海岬科技相关，\texttt{E-204}的属性与东岭大学相关。对“林舟介绍了海岬科技的新系统”，规则分两步走：先按名字把两个实体都取为候选，再借助额外证据\emph{消歧}——句中出现了海岬科技、或文档带有指向某作者主页的标识、或来源本身就与某机构相关，据此判断这个提及应当对应哪位作者。要注意，句中出现海岬科技只是一条\emph{待检查的线索}，“名字相同”这一点本身完全没有确定身份，不能跳过证据直接拍板。

如果上下文极少，只有一句“林舟发表了报告”，没有任何能区分两位的线索，规则就应当老实地保留候选集合$\{\texttt{E-017},\texttt{E-204}\}$，而不是任意挑一个——猜错的代价远大于保留未知。还要区分两种“拿不到唯一实体”的不同情况：库里\emph{一个候选都没有}，说明知识库可能缺这个实体，需要补充实体或检索线索；库里\emph{有多个候选却无法区分}，说明实体存在但证据不足，需要补充能辨别身份的上下文。两种情况的补救方向不同，应分别保存。

除了用上下文消歧，还有几种常见的链接手段，各自的可靠性边界不同。如果提及本身带有序号、主页链接这类\emph{唯一标识}，按标识查询几乎可以直接确定实体，是最可靠的情形。别名表把简称、旧名都作为查询入口，能召回更多写法，但别名往往有时间性（机构曾用名、人物婚前姓），需要记录名称在什么时间段有效。模糊字符串匹配放宽了对拼写完全一致的要求，能容忍错别字和写法变体，但放得越宽，越可能把“名字相近、其实不同”的两个实体混为一谈。模型可以利用更广的上下文辅助判断该选哪个候选，但它承担的仍然是“在候选中选择关联对象”这一步，不能替代对最终链接的核验。

把实体链接拆成两个阶段在排错时特别有用，这也是实体链接研究中的标准划分\scite{knowledge-acquisition-milne2008}：先做\emph{候选生成}（按名字、别名等从库里检索出可能的实体集合），再做\emph{基于上下文的候选选择}（从中挑一个或保留多个）。这样拆开后，错误就能对号入座：最终链接错了，可能是因为正确实体压根没进入候选集（召回失败），也可能是它在候选集里、却被选错了（排序失败），两者的改进方法完全不同。

链接到实体只是手段，查到属性之后还要再做一次“属性—任务”对照，否则容易把看似相关实则不对口的属性用错。例如实体\texttt{E-017}的职业属性写着“软件工程师”，但任务要回答的是“文档所述那个时间点上他担任什么职务”；“软件工程师”是一个相对静态的职业身份，并不等同于他在特定时间、特定机构的具体任职关系，二者不能互相替代。

这一步要能区分三类不同来源的错误：其一是\emph{事实本身有误}（库里的属性就是错的）；其二是\emph{链接错误}（属性是对的，但提及被接到了错误的实体上）；其三是\emph{属性不适用}（实体对、属性也对，但这条属性不是当前任务或当前时点所问的那个）。为了事后能区分三者，规则应保存实体库版本、属性的生效时间、当初的候选列表、消歧所用证据以及最终选择，形成一条完整可回溯的链接链。

实体词表也需要持续维护。新别名可以由专家补充、从历史匹配样本中归纳，或采纳模型的建议，但每加入一个别名都要检查两件事：它会不会其实是另一个实体的名字（导致两个实体被错误合并），以及它是否只在特定时间段有效。改动别名后，还应回看旧的标注结果，确认这次改动会影响哪些历史链接、是否需要重算。

评估实体链接质量时，要按实体分组、并分别盯住前述两个阶段，因为不同实体的“知名度”差异很大。对每个目标实体组，分别检查\emph{候选召回}（正确实体到底有没有进入候选集）和\emph{选择错误}（在候选里是否选错了身份）。一个重要的反直觉点是：关于某个实体的材料多，既不保证它的召回率更高，也不保证消歧更准——材料多有时反而带来更多歧义别名。因此，在资料充分的热门实体上测得的好结果，不能直接外推到资料稀少的冷门实体，后者往往才是错误集中的地方。

\subsection{基于概念关联的规则利用}
\label{sec:rule-based-knowledge-acquisition-concept-association}

实体匹配解决“这个名字指谁”，但新闻里的“首席架构师”还面临另一个独立问题：它在任务规定的岗位分类里属于哪一类？这与确定“林舟”是哪位作者是\emph{不同}的问题——判断头衔的岗位类别，并不需要先确定戴这个头衔的人的唯一身份，哪怕有两位同名的林舟，“首席架构师”这个头衔的归类方式也是一样的。

\term{概念关联}（concept association）处理的就是这种“词义到类别”的对应：把文本中的词语或短语对应到词典、本体或某个分类体系中的概念标识，然后依据该体系明确规定的上下位（从属）关系，推出类别或缩小候选范围。这里有一个清晰的分工：\emph{文字}只负责在映射表里找到对应概念，真正的\emph{推导}则依据概念标识与概念之间在体系中的关系完成，而不是再靠字面猜测。

沿概念层级推导时，\textbf{方向}至关重要，最容易出错的就是把上下位关系反着用。假设教学岗位体系把“首席架构师”对应到细概念\texttt{role:\allowbreak{}technical-\allowbreak{}lead}，又规定该概念从属于上位类\texttt{role:\allowbreak{}management-\allowbreak{}position}。那么当文本确认了“首席架构师”这个细类后，规则可以\emph{沿从属关系向上}推出“管理岗位”这一上位类别，这个推导是成立的。

但反过来不成立：只知道某人属于“管理岗位”，绝不能反推出他就是“首席架构师”，因为管理岗位这个上位类之下还包含技术总监、产品负责人等许多其他岗位——上位类信息不足以唯一确定下位类。这个方向差别直接决定了规则能输出什么：做\emph{粗类别}判断的任务，上位概念就够用，可以直接给类；做\emph{细类别}判断的任务，若缺乏进一步证据，只能保留下位候选集合、标为无法唯一确定，不能硬选一个。还要注意，这里的归类结论依据的是\emph{指定的这套教学体系}，换成另一套岗位划分方式结果可能不同，不能把它当成对所有岗位体系的普遍断言。

建立“词语到概念”映射时，不能把词与词之间的各种关系混为一谈，因为它们允许的操作不同。在当前用法下确属同义的词，可以指向同一个概念；缩写词需要结合领域和上下文才能确定展开成哪个概念（同一个缩写在不同领域可能对应完全不同的词组）。尤其要防止两种过度合并：\emph{近义词}只是含义相近，在任务要求精确区分时不足以认定为同一概念；\emph{上下位词}（如“工程师”与“首席架构师”）之间存在层级方向，更不能合并成一个概念而丢掉从属关系。

WordNet是这类资源的经典范例：它先把严格意义上的同义词组织成一个个“同义词集”（synset），每个同义词集对应一个概念，再用上位、下位等语义关系把这些概念连接起来\scite{knowledge-acquisition-miller1995}。这种“词面集合”与“概念关系”分开记录的做法值得借鉴——它迫使使用者区分“两个词能不能互换”和“两个概念谁从属于谁”这两类问题。

为了让归类结果可追查，三类知识应当\emph{分开维护}，因为它们回答的是链路上不同的问题：概念映射表确定“这段文本指向哪个概念”，本体关系确定“概念与概念之间是什么关系”，推导规则确定“哪些关系可以、以及如何转成任务标签”。混在一张表里，出错时就无法判断是映射错了、关系定义错了、还是推导规则用错了。映射候选同样可以由人工整理、从词语使用模式归纳或由模型建议，但无论来源如何，都要分别核验“文本表达的含义是否真对应该概念”和“所用的体系关系是否正确”。

使用层级时还要警惕\emph{属性继承}被滥用。沿上下位关系能继承\emph{类别成员资格}（首席架构师是管理岗位），但并非该概念的一切属性都能无条件向下或向上传递。某个具体岗位的薪酬区间、任期长度、资格要求，未必适用于上位类下的所有岗位；反过来，上位类的某些泛化属性，下位的特殊岗位也可能有例外。能不能继承某条属性，要回到该属性的语义本身去核对，不能因为“是子类”就默认一切属性都继承。

概念关联还有一个不可回避的多体系问题：同一个岗位词，在“按职责划分”“按组织级别划分”“按行业标准划分”的不同体系里，可能落在完全不同的分类位置。因此规则必须记录它当前采用的是\emph{哪一套体系、哪个版本}；当体系发生概念合并、拆分或层级调整时，才能沿“映射到推导”的链条，反查出哪些旧结果需要重算。

最后同样要守住“未关联并非负类”的边界：当前体系里找不到某个词的概念映射，只说明\emph{规则没能建立关联}（可能是词典没收这个新词），并不能证明这个样本就不属于目标类别。这与模式匹配中“未触发不等于负例”是同一个道理。至于使用了多套体系、需要把不同来源的概念彼此对齐时该怎么处理，则留到第~\ref{chap:knowledge-fusion}章讨论。

\subsection{基于关系匹配的规则利用}
\label{sec:rule-based-knowledge-acquisition-relation-matching}

实体匹配确定了“谁是谁”、概念关联确定了“属于哪类”，但若任务要标注的是“林舟\emph{任职于}海岬科技”，只查单个实体的属性或岗位类别仍然不够——这里要判断的是\emph{两个实体之间}的关系。\term{关系匹配}（relation matching）把样本中的一对实体与知识库中的关系事实对应起来，再用关联到的事实产生关系标签，或成组地提供监督信号。

从单实体转到关系，检查对象发生了实质变化：现在的基本单位是一个\emph{有次序的实体对}（谁在前、谁在后不能颠倒），而且必须同时确定关系的定义和它的有效时间。例如三元组$(\text{林舟},\text{任职于},\text{海岬科技})$是有方向的，不能换成$(\text{海岬科技},\text{任职于},\text{林舟})$（方向反了，关系就不成立）；它也有时间性，不能不加区分地套用到该事实生效之前或失效之后的所有文本上。

关系标注需要大量“这句话是否表达了某关系”的逐句标签，全靠人工标注极其昂贵。一个能批量、自动产生训练候选的想法应运而生：既然知识库已经记录了“林舟任职于海岬科技”，那么凡是\emph{同时提到}这两个实体的句子，很可能都在表达这项关系——于是直接用知识库事实给这些共现句子打上关系标签。这种\emph{借助已有知识库事实、自动为提及它的文本提供标签}的办法，称为\term{远程监督}（distant supervision）。Mintz等人的经典设定正是把知识库关系与文本中共同出现的实体对对齐，据此自动构造训练实例\scite{knowledge-acquisition-mintz2009}。

必须立刻点明这个假设的强度：远程监督依据的是一条启发式——“若两个实体存在某种关系，则同时提到它们的句子倾向于表达该关系”。它\emph{没有逐句证明}每一个共现句子真的在表达那项事实。共现只是关系表达的必要线索，远非充分条件，下一段的对照就会暴露由此产生的错误标签。理解这一点，是正确使用远程监督的前提：它能极大降低标注成本，但代价是标签里混入噪声。

用一个严格控制的教学对照看清噪声从何而来。继续使用构造事实“林舟任职于海岬科技”，并把目标关系限定为“文本所述时间已经生效的任职状态”；假设句甲中的任命已经生效、所述时点仍在该事实的有效期内。在这些明确前提下比较两句话：句甲“海岬科技任命林舟为首席架构师”确实支持目标任职关系；句乙“林舟在论坛上批评了海岬科技的发布流程”虽然也同时提到了林舟和海岬科技，却只表达了“批评”这一动作，与任职毫无关系。

如果机械地执行远程监督、把事实复制给所有共现句子，句乙就会被错误地贴上一个“任职关系”的正标签，而这个标签在文本中毫无依据——这正是远程监督噪声的典型来源。需要说明，这里“句甲支持、句乙不支持”是教学中我们\emph{已知}的语义对照，用来帮助读者理解问题，并不是机器仅凭共现就能作出的判断；真实系统面对成百万条句子时并不知道哪一句是“句乙”。图~\ref{fig:rule-based-knowledge-acquisition-relation-path}中，两句经过实体消歧和关系对齐后会被归入同一个“提及组”；图中表示语义对照的点线只是教学标注，不属于自动判定步骤，构造提及组时也不要求预先知道组内哪一句真正表达了关系。

\begin{figure*}[htbp]
  \centering
  \begin{tikzpicture}[
  x=1cm,
  y=1cm,
  font=\figurefont\footnotesize,
  source/.style={draw=FigureInput,line width=0.9pt,fill=FigureInput!9,minimum height=10mm,text width=32mm,align=left,inner sep=1.5mm},
  fact/.style={draw=FigureModel,line width=1pt,fill=FigureModel!10,minimum height=9mm,text width=53mm,align=center,inner sep=1.5mm},
  process/.style={draw=BookMuted,line width=0.8pt,fill=BookPaper,minimum height=10mm,text width=13mm,align=center,inner sep=1.5mm},
  comparison/.style={draw=BookMuted,line width=0.8pt,dotted,fill=white,minimum height=15mm,text width=36mm,align=center,inner sep=1.5mm},
  group/.style={draw=FigureOutput,line width=1.1pt,fill=FigureOutput!7,minimum height=37mm,text width=36mm,align=center,inner sep=1.5mm},
  flow/.style={knowledge arrow,draw=BookInk,line width=1.1pt,rounded corners=1mm},
  evidence/.style={knowledge arrow,draw=FigureModel,line width=0.9pt,dashed,rounded corners=1mm},
  contrast/.style={draw=BookMuted,line width=0.8pt,dotted},
  edge label/.style={text=BookMuted,inner sep=1pt,font=\figurefont\scriptsize}
]
  \fill[white] (0,0) rectangle (16.7,7);

  \node[source] (sent-a) at (1.8,5.35) {句甲：海岬科技任命林舟为首席架构师。};
  \node[source] (sent-b) at (1.8,2.25) {句乙：林舟在论坛上批评了海岬科技的发布流程。};

  \node[process] (entity-a) at (4.85,5.35) {实体消歧};
  \node[process] (entity-b) at (4.85,2.25) {实体消歧};
  \node[process] (mapping-a) at (7.1,5.35) {关系对齐};
  \node[process] (mapping-b) at (7.1,2.25) {关系对齐};

  \node[fact] (kb) at (7.35,3.8) {已有构造事实\\$(\text{林舟/E-017},\text{任职于},\text{海岬科技/E-003})$};
  \node[comparison] (positive) at (10.3,5.35) {教学语义对照\\句甲表达关系\\可作句子级候选};
  \node[comparison] (cooccur) at (10.3,2.25) {教学语义对照\\句乙仅共同提及\\不应给句子正标签};
  \node[group] (bag) at (14.55,3.8) {提及组级候选\\句甲、句乙组成一组\\由实体关系库事实关联关系\\假设至少一条表达关系};

  \draw[flow] (sent-a) -- (entity-a);
  \draw[flow] (entity-a) -- (mapping-a);
  \draw[contrast] (mapping-a.east) -- (positive.west);
  \draw[flow] (mapping-a.north) -- (7.1,6.55)
    -- node[edge label,above] {按实体对归组} (14.55,6.55) -- (bag.north);

  \draw[flow] (sent-b) -- (entity-b);
  \draw[flow] (entity-b) -- (mapping-b);
  \draw[contrast] (mapping-b.east) -- (cooccur.west);
  \draw[flow] (mapping-b.south) -- (7.1,1.05)
    -- (14.55,1.05) -- (bag.south);

  \draw[evidence] (kb.north -| mapping-a.south) -- (mapping-a.south);
  \draw[evidence] (kb.south -| mapping-b.north) -- (mapping-b.north);
  \draw[evidence] (kb.east) -- (bag.west);

  \draw[flow] (0.35,0.35) -- (1.25,0.35);
  \node[anchor=west,text=BookMuted,font=\figurefont\scriptsize] at (1.35,0.35) {处理步骤};
  \draw[evidence] (4.1,0.35) -- (5.0,0.35);
  \node[anchor=west,text=BookMuted,font=\figurefont\scriptsize] at (5.1,0.35) {事实关联};
  \draw[contrast] (7.85,0.35) -- (8.75,0.35);
  \node[anchor=west,text=BookMuted,font=\figurefont\scriptsize] at (8.85,0.35) {已知语义对照，非自动判定};
\end{tikzpicture}

  \caption{从已有事实到文本候选关系的关联路径。实体和关系均为构造例子，并假设句甲的任命已生效、所述时点在事实有效期内；实线箭头表示处理步骤，虚线箭头表示实体关系库事实的证据关联，点线引出教学中已知的语义对照，不表示自动判定。实体消歧确认两句中的提及指向同一实体对；关系对齐检查定义、方向与时间。在上述前提下，句甲支持任职关系，句乙仅共同提及，但两句都可按实体对进入同一提及组。组级候选依赖至少一次表达的假设，不要求预先知道哪句支持关系，也不能据此给每句确定的正标签。}
  \label{fig:rule-based-knowledge-acquisition-relation-path}
\end{figure*}

要缓解上述噪声，需要先把整条关联路径上的几个环节走扎实：确认实体身份与关系方向，核对知识库关系与任务关系的定义是否对齐，最后还要明确一个关键问题——\emph{证据应当出现在什么范围}：是要求就在当前这一句里表达，还是放宽到整段话、乃至同一实体对的多条提及中只要有一次表达即可？证据范围一旦改变，被监督的对象（单句 vs. 实体对）也随之改变。

远程监督的原始设定近似要求“每个共现提及都表达该关系”，过强而不符合实际。Riedel等人把它放松为\emph{至少一次表达}的假设：分别建模两个层次——“这个实体对\emph{总体上}是否存在该关系”和“\emph{某一条具体提及}是否表达该关系”，只假设一个存在关系的实体对，其多条提及中\emph{至少有一条}真正表达了关系，而不要求条条都表达\scite{knowledge-acquisition-riedel2010}。这样一来，像句乙那样的噪声提及就可以被容忍：它不必为正，只要同组里存在像句甲那样真正表达关系的提及即可。

在“至少一次”假设下，数据的组织方式随之改变：把同一实体对的多条提及打包成一个\emph{提及组}（bag），赋予组一个\emph{组级候选}标签，表示“组内至少有一条提及可能支持该关系”，而不再给组内每一句话都打正标签。模型学习时可以先对每一句“是否表达关系”分别打分，再按某种方式（如取最大分）聚合成组级预测，训练时只把组级预测与组级标签比较。第~\ref{chap:knowledge-acquisition-model-based}章的式~\eqref{eq:knowledge-acquisition-model-bag-prediction}给出了这种多实例预测的具体计算和所依赖的假设。

要厘清本章规则在其中的角色：规则负责产出提及组、实体标识和关联到的知识库事实，即把原料和组结构准备好；但“组内究竟哪一句构成证据”仍是一个需要\emph{提及级推断}的学习问题，规则并不回答。因此不能用教学里我们人为已知的“句甲支持、句乙不支持”去代替模型本应输出的预测——那是讲解用的答案，不是自动流程的产物。

把整条路径（实体消歧、关系对齐、证据范围判断）的每一步都记录下来，出错时才能区分性质不同的故障：消歧环节出错会把提及关联到无关实体；使用了过时事实会把关系套到错误的时间；而文本本身没表达关系，则是证据范围条件未被满足。

还有一种常见情况：知识库里\emph{根本没有}查到该实体对的关系。这时应采取什么态度至关重要。本章采用\term{开放世界假设}（open-world assumption）：承认知识库可能是不完整的，“没有查到这条关系”只说明库中未记录，\emph{不等于}“已证明该关系不存在”——也许只是还没收录。因此规则在查不到时应保留未知或候选关系，而不是草率地返回负标签；只有当存在明确的排除性证据、且它正好符合任务对负类的定义时，才能给出负例。这与全章反复强调的“未触发/未查到不等于负类”一脉相承。

开放世界假设带来大量“只有正例已知、其余未知”的实体对，自然引出一个问题：训练关系分类器时，是否\emph{必须}先把这些未知对象猜成负例？把未知一律当负，显然会把其中真正的正关系误标，扭曲训练。

其实不必如此。在一定抽样条件下，可以直接用“一批确认的正例”加“一批未标注样本（其中正负混杂）”来学习，这就是正例—未标注学习（PU learning）。其条件是：正例样本能够代表目标正类的分布、未标注样本能够代表目标总体分布、并且正类比例已知或可估计；满足这些条件时，可以从两批样本的平均损失\emph{重建}出完整的分类风险，而无须逐条把未知改写成负类。具体构造见第~\ref{sec:ka-pu-learning}节。需要强调这些条件并不免费：如果知识库偏向于只收录显著、知名的关系，那么抽中的正例往往是“显著关系”的代表、而非全部正关系的代表，代表性假设就不能默认成立；开放世界假设本身只说明“未查到并不等于不存在”，它并不自动提供PU学习所需的抽样保证。

即使上述各步都能正确执行，最终候选的\emph{分布}仍可能失真，需要单独检查两类量：每个关系类别的候选总量，以及每个实体对的提及组大小。读这些数字时要避免错误归因：两个关系类别的候选总数不同，可能只是因为它们涉及的实体对数量本来就不同，不能据此推断“某类中单个实体对被提及得更多”——总量与组内密度是两回事。

还要警惕重复带来的虚高：如果同一实体对在新闻里被反复共现，或相关文本被多家媒体转载，在没有去重的情况下，这些近乎重复的提及都会被计入该实体对的提及数，使组显得很大、证据很“充分”，但信息其实是重复的。这类分布差异仅靠“规则执行成功”无法发现，因此还要保存知识库版本、事实来源、实体对、组内提及列表以及去重标识，供后续的融合与可靠性检查追溯。

\section{基于约束推理的规则利用}
\label{sec:rule-based-knowledge-acquisition-constraint}

前两类规则都是在\emph{直接产生}答案：模式匹配说“这个样本像正类”，知识关联说“该实体有某属性”。但还有一类常见情形，单个答案看起来都合理，放在一起却自相矛盾。比如对同一岗位，细类判断说“他是首席架构师（为真）”，上位类判断却说“他不属于管理岗位（为假）”，而按岗位体系首席架构师必属管理岗——两个结论不可能同时成立。

处理这种矛盾，最自然的办法不是再分别重判一次，而是退一步规定：\emph{这些答案之间必须满足什么关系}。这就引出第三类依据。\term{约束推理}（constraint-based inference）先把每个未知标签所有可能的取值列为候选，然后用逻辑、结构或领域知识给出的条件，\emph{排除掉那些不允许的组合}；它通常不直接指定唯一答案，而是缩小可行答案的范围。排除之后若只剩一个候选，答案随之确定；若仍剩多个，可以把这一整个范围作为结果保留，也可以再引入一类“允许违反、但违反要付出代价”的软偏好，在多个可行候选中选取代价最小的。

要准确说出“约束到底排除了什么”，先要划定一个出发点：在不考虑任何约束时，任务\emph{理论上允许}哪些完整答案？把所有这些完整答案收集成集合$\bm{Y}$，其中每个元素$\bm{y}$是对一组未知标签或数值的一次\emph{完整赋值}（即给每个未知量都指定一个取值）；$\bm{Y}$由任务规定的取值范围确定，比如每个二值标签可取0或1，则$n$个标签共有$2^n$种赋值。

然后再施加必须无条件满足的条件。设有$m$条这样的硬约束$c_1,\ldots,c_m$，每条约束对完整赋值$\bm{y}$返回真或假。\emph{同时满足全部硬约束}的那些赋值构成可行集合：
\begin{equation}
  \bm{F}
  =
  \{\bm{y}\in\bm{Y}:c_k(\bm{y})=\mathsf{T},\ k=1,\ldots,m\}.
  \label{eq:rule-based-knowledge-acquisition-feasible-set}
\end{equation}
式中要求每一条硬约束同时为真，因此$\bm{F}$回答的问题是“哪些\emph{完整答案仍然被允许}”，请特别注意它\emph{不}回答“哪个答案更好、更受偏好”——可行与否是一个是非判断，偏好排序则是另一回事（将由软约束处理）。

在约束确实能对所有候选求值的前提下，可行集合$\bm{F}$会出现三种情况，必须区别对待，因为它们传递的信息完全不同。其一，$\bm{F}$只有唯一元素：当前前提恰好把答案完全确定，这个唯一赋值可作为结果。其二，$\bm{F}$有多个元素：依据只足以把范围缩小，但不足以唯一确定，此时应保留整个可行范围而不是任取一个冒充确定答案。其三，$\bm{F}$为空集：没有任何赋值能同时满足所有条件，这是一个有价值的信号——说明已知事实、约束规则或观测值之间\emph{存在冲突}（例如数据违反了物理守恒），需要排查，而不是随便给个答案。还要区分“无解”与“无法判断”：若只是缺少某些前提、导致约束暂时无法求值，应保留为无法确定，而不能草率宣告无解。

把这些映射回样本标注状态。当任务\emph{允许返回范围}时，一个非空的可行集合本身就可以作为内容提交：来源动作为作答、语义状态为已给出，但要意识到其中单个值可能仍未知。而当无法判断（缺前提）或已证明无解（$\bm{F}$为空）时，本节都规定\emph{不}提交任务赋值：来源动作为弃权、内容为空、语义为未知；但两者要分别留痕——缺前提时保存“缺的是哪个前提”，无解时另存无解诊断和相互冲突的依据，不能把“算不出来”和“算出来发现矛盾”混成一类。与前面一样，上述执行状态都以“相应的检查或求解确实跑完”为成功前提；程序中断记为执行失败。最后，求解顺利完成只在逻辑上得到了可行集，验证状态并不会因此自动变成“已核验”——满足约束不等于符合事实。还要记住触发与求解的区别：条件检查触发了求解，绝不代表求解已经给出任务内容。

\subsection{基于逻辑约束的规则利用}
\label{sec:rule-based-knowledge-acquisition-logical-constraint}

最基本的约束作用于若干真假命题之间的逻辑关系。要排除“细类成立、上位类却不成立”这种矛盾，可以把两项类别判断分别写成命题，再规定它们之间的\emph{蕴含}关系。一般地，\term{逻辑约束}（logical constraint）用蕴含（“若$p$则$q$”）、互斥（“$p$与$q$不能同时成立”）、等价（“$p$与$q$同真同假”）等关系，限制若干命题允许怎样\emph{共同}取值。

逻辑约束有两种强弱不同的用法。当作\emph{硬约束}时，它代表必须遵守的规则，任何违反该关系的候选组合都直接从可行集合中剔除，没有商量余地。但有时我们想表达的并不是“绝对禁止”，而是“最好如此、违反了会差一些”的\emph{偏好}，这时可改用\term{软约束}（soft constraint）：违反它的候选仍然允许保留，但要为此付出一个事先指定的代价。由于违反软约束不会使候选不可行，最终选谁就必须结合代价的权重或其他比较依据，而不能只看是否违反。

用最简单的两个命题把逻辑约束算一遍。设二值命题$p$表示“属于技术管理岗位”，$q$表示“属于管理岗位”，教学体系要求$p\Rightarrow q$（技术管理岗位必然是管理岗位）。两个二值命题共有$2^2=4$种赋值组合：$(0,0),(0,1),(1,0),(1,1)$。逐一检验，只有$(p,q)=(1,0)$违反蕴含（是技术管理岗却不是管理岗），被硬约束排除；其余三种仍在可行集合中。这一步只是在做逻辑筛选，并没有判断现实中$p,q$到底取什么值。

如果还想在三个可行候选中进一步挑一个，就需要引入偏好：人为给每个候选设置一个软约束代价$L(\bm{y})$，然后按代价从小到大比较。图~\ref{fig:rule-based-knowledge-acquisition-feasible-set}给了一组教学代价，其中三个可行项里最小者恰好唯一。但这里有一个必须向读者讲透的要点：\textbf{代价最小只代表它最符合“我们当前设定的偏好”，完全不证明它就是真实答案}。图的下方面板则演示了另一种情形：当观测之间本身矛盾（比如同时观测到$p=1$、$q=0$，又坚持$p\Rightarrow q$），没有任何赋值能满足全部条件，可行集合变为空——这提示数据或约束冲突，而不是某个候选“代价高”。

\begin{figure*}[htbp]
  \centering
  \begin{tikzpicture}[
  x=1cm,
  y=1cm,
  font=\figurefont\footnotesize,
  candidate/.style={draw=BookMuted,line width=0.8pt,fill=BookPaper,minimum width=30mm,minimum height=13mm,align=center,inner sep=1.5mm},
  selected/.style={candidate,draw=FigureModel,line width=1.2pt,fill=FigureModel!10},
  premise/.style={draw=FigureInput,line width=0.9pt,fill=FigureInput!9,minimum width=28mm,minimum height=10mm,align=center,inner sep=1.5mm},
  constraint/.style={draw=BookAmber,line width=1pt,fill=BookAmber!9,minimum width=31mm,minimum height=10mm,align=center,inner sep=1.5mm},
  combine/.style={circle,draw=BookInk,line width=0.9pt,fill=white,minimum size=7mm,inner sep=0pt},
  empty/.style={draw=BookAmber,line width=1.2pt,fill=BookAmber!9,minimum width=32mm,minimum height=11mm,align=center},
  flow/.style={knowledge arrow,draw=BookInk,line width=1.1pt,rounded corners=1mm},
  soft/.style={text=BookMuted,font=\figurefont\scriptsize}
]
  \fill[white] (0,-1.15) rectangle (16.6,6.5);
  \node[anchor=north west,font=\figurefont\bfseries] at (0,6.25) {(a) 硬约束筛选与软偏好};
  \node[constraint,minimum width=44mm] (hard) at (8.3,5.75) {硬约束：$p\Rightarrow q$};

  \node[candidate] (c00) at (1.8,4.25) {$(p,q)=(0,0)$\\{\scriptsize 可行；$L=1.0$}};
  \node[candidate] (c01) at (5.65,4.25) {$(p,q)=(0,1)$\\{\scriptsize 可行；$L=0.6$}};
  \node[candidate,draw=BookAmber,fill=BookAmber!7] (c10) at (9.5,4.25) {$(p,q)=(1,0)$\\{\scriptsize 违反硬约束}};
  \node[selected] (c11) at (13.35,4.25) {$(p,q)=(1,1)$\\{\scriptsize 可行；$L=0.2$（最小）}};
  \draw[BookAmber,line width=1.2pt] ([xshift=-10mm,yshift=-4mm]c10.center) -- ([xshift=10mm,yshift=4mm]c10.center);
  \draw[BookAmber,line width=1.2pt] ([xshift=-10mm,yshift=4mm]c10.center) -- ([xshift=10mm,yshift=-4mm]c10.center);
  \node[soft,anchor=north] at (8.3,3.42) {硬约束留下三个答案；软代价只负责排序};

  \draw[BookRule,line width=0.8pt] (0,2.95) -- (16.6,2.95);
  \node[anchor=north west,font=\figurefont\bfseries] at (0,2.75) {(b) 新增互相矛盾的已知条件};
  \node[premise] (p-one) at (2.2,1.65) {观测：$p=1$};
  \node[premise] (q-zero) at (7.2,1.65) {观测：$q=0$};
  \node[constraint] (imply) at (12.2,1.65) {硬约束：$p\Rightarrow q$};
  \node[combine] (and) at (7.2,-0.35) {$\land$};
  \node[empty] (none) at (12.4,-0.35) {无可行解\\报告冲突};
  \draw[flow] (p-one.south) |- (and.west);
  \draw[flow] (q-zero.south) -- (and.north);
  \draw[flow] (imply.south) -- ++(0,-0.5) -| (and.45);
  \draw[flow] (and.east) -- (none.west);
\end{tikzpicture}

  \caption{约束缩小可行答案并产生冲突的教学示意。上面板从两个二值命题的四种赋值出发，硬约束$p\Rightarrow q$排除$(1,0)$，软约束代价再在三个可行候选中给出唯一最小值；这些代价是人为设置，不是统计估计。下面板同时加入观测$p=1$、观测$q=0$和同一蕴含后，可行集合变为空集。逻辑可行、偏好最小和真实正确是三种不同判断。}
  \label{fig:rule-based-knowledge-acquisition-feasible-set}
\end{figure*}

实际偏好往往不止一条，需要一个可计算的量来汇总多条软约束的共同作用。令第$j$条软约束对候选$\bm{y}$的违反代价为$\ell_j(\bm{y})\geq0$：取零表示完全不违反，取正值表示有违反，值越大违反越严重；再给每条软约束配一个权重$w_j\geq0$，规定这一项偏好在总比较中占多大分量。

一个关键安排是：\textbf{软约束只在硬约束筛过的可行集合$\bm{F}$内部起作用}。也就是说，先无条件满足全部硬约束，再对选定的有限条软约束求加权代价之和，总数越小表示越受偏好。绝不允许为了降低软代价而选择违反硬约束的候选——硬约束是入场资格，软约束只是入场后的排序依据。当可行集非空、且最小值确实能取到时，选择总加权代价最小的候选，得到
\begin{equation}
  \widehat{\bm{y}}
  \in
  \argmin_{\bm{y}\in\bm{F}}
  \sum_j w_j\ell_j(\bm{y}).
  \label{eq:rule-based-knowledge-acquisition-soft-constraint}
\end{equation}
式中$\bm{y}\in\bm{F}$这一限制保证了被比较的候选都没有违反硬约束；目标值更小，才表示它在“当前给定的权重”下更受偏好——换一套权重，最优候选可能改变。$\argmin$会保留\emph{所有}达到最小值的候选而不强行挑一个，因此若最小值由多个候选并列取得，答案实际上仍未唯一确定，此时应如实报告多解。存在性也要分清：在有限的非空候选集合上、代价又有限，最小值一定存在；但若是连续取值的优化问题，最小值未必能取到（可能只是下确界），需要另外确认前述存在条件。

这种“硬约束定边界、软代价做排序”的框架，可以通过整数规划或线性规划来求解，从而把多个局部模型的预测与任务层面的逻辑条件融合起来；Roth和Yih在自然语言任务中给出了有代表性的形式化\scite{knowledge-acquisition-roth2004}。还可以再往前走一步：若希望把这套加权逻辑\emph{解释成一个概率模型}（让权重对应“某条约束成立的可能性”），就必须明确权重的统计含义，并确认公式代入具体对象后有足够的实例支撑；Markov逻辑网络为此提供了一种把加权逻辑规则概率化的实现方式\scite{pgm-energy-models-richardson2006}。但要注意，随手设定的权重并不天然具备概率解释。

把知识写成逻辑约束时，最容易犯的错误是丢掉关系的\emph{方向}和\emph{适用对象}，这与概念关联中上下位不可逆是同一个警示。由“首席架构师属于技术负责人”只能得到有方向的蕴含$\text{首席架构师}\Rightarrow\text{技术负责人}$，绝不能反过来说所有技术负责人都是首席架构师。此外，每条约束往往只在特定对象范围内成立：针对“正式岗位”设定的互斥条件（一人不能同时任一岗）未必适用于允许兼任的临时角色；把约束的适用范围悄悄扩大，会错误地砍掉合法候选。

因此，无论约束来自专家编写、从样本归纳还是模型建议，都要逐项核对四件事：蕴含方向有没有写反、量词（对“所有”还是“某些”对象成立）是否正确、作用对象是否恰当、以及有没有被忽略的例外情况。

为了让约束结果可维护、可在规则更新后复查，一次推理应留下完整的“决策说明”：用了哪些前提、据此排除了哪些候选、每条软约束分别贡献了多少代价。这样当某条约束或权重调整后，才能沿说明重新计算并看出结果为何改变。成本也要如实计入：约束求解本身的计算量、软权重如何设定、出现冲突时如何解释，都不是免费的。

最后划清本章与下一章的边界：这里讨论的约束，作用是\emph{从既定前提出发去产生或限制新的候选}；如果约束主要被用来\emph{综合多个已经存在的来源}给出的答案（比如调和多个模型、多种标注的分歧），那属于第~\ref{chap:knowledge-fusion}章的融合方法，不应在规则推理层提前完成。

\subsection{基于结构约束的规则利用}
\label{sec:rule-based-knowledge-acquisition-structural-constraint}

逻辑约束处理的是抽象命题之间的关系，但很多任务的标签带有\emph{空间或序列结构}：一个实体名可能横跨好几个词，每个词各自选了标签之后，还必须保证这些标签能拼出一个合法的片段。

\term{结构约束}（structural constraint）利用样本在序列、树、图或成组形式中的组织关系，来限制标签\emph{之间}如何组合。它与一般逻辑约束的区别在于：除了标签的取值关系，还必须先明确“哪些位置（或对象）彼此相邻、谁包含谁、哪些属于同一组”——也就是约束作用的那张“结构”从哪来。这张结构并非凭空假设，它可以直接来自任务格式（如固定的字段顺序）、标注规范（如片段如何界定）或样本字段，也可以由实体匹配等前置步骤构造出来；来源不同，可靠性也不同。

要让一串逐词标签表达“片段从哪开始、到哪结束”，需要一套边界编码约定。本节采用BIO体系中的IOB2约定，三个标记分别是beginning（开始）、inside（内部）和outside（外部）的缩写：对任意实体类型X，一个X类型片段的第一个词标为\texttt{B-X}，片段内部后续的词标为\texttt{I-X}，不属于任何片段的词标为\texttt{O}。

这套约定立刻导出一条结构约束：\texttt{I-X}只能出现在\texttt{B-X}或另一个\texttt{I-X}之后（即延续一个已经开始的同类型片段），它\emph{不能凭空出现}在句首或一个\texttt{O}之后——否则就成了一个“没有开头的片段”，是非法组合。这种把片段结构编码进逐词标签、再检查局部转移合法性的做法，源自早期的文本分块工作；需要提醒的是，历史上还有一种IOB1约定，它只在两个同类片段紧临时才用\texttt{B}标记后一个片段的开头，其余片段内部用\texttt{I}开头，与本节“每个片段都以\texttt{B}开始”的IOB2约定不同，阅读文献时不要混淆\scite{knowledge-acquisition-ramshaw1995}。

用一个四词教学序列把约束的作用和局限都看清楚。对“林舟~加入~海岬~科技”，先由局部规则给每个位置提供候选标签，其中\texttt{PER}、\texttt{ORG}分别表示人员和组织。假设逐点独立选择后得到一串标签
\[
  \texttt{B-PER},\ \texttt{I-ORG},\ \texttt{B-ORG},\ \texttt{I-ORG}.
\]
检查相邻转移就会发现非法之处：第二个位置的\texttt{I-ORG}想延续一个组织片段，但它前面的位置是\texttt{B-PER}（人员片段的开头），并不存在一个已开始的组织片段可供延续，违反了“I必须承接同类型B/I”的规则，因此整串组合非法，必须修改。图~\ref{fig:rule-based-knowledge-acquisition-sequence-constraint}给出的另一串\texttt{B-PER, O, B-ORG, I-ORG}则满足所有转移条件，是一个合法序列。

但这个例子同时暴露了结构约束的\emph{局限}：合法不等于正确。把“海岬科技”两个字全都标成\texttt{O}（序列变成\texttt{B-PER, O, O, O}）其实也完全满足转移规则、照样合法，可它显然漏掉了真正的组织提及。结构约束只能排除“拼不成合法片段”的组合，却无法判断众多合法组合中哪一个在语义上正确；要选出真正的任务答案，还必须借助词语、上下文或其他来源提供的\emph{语义}证据。约束负责“剔除不可能”，语义负责“在可能中选对”。

\begin{figure*}[htbp]
  \centering
  \begin{tikzpicture}[
  x=1cm,
  y=1cm,
  font=\figurefont\footnotesize,
  token/.style={draw=FigureInput,line width=0.9pt,fill=FigureInput!9,minimum width=30mm,minimum height=8mm,align=center},
  tag/.style={draw=FigureModel,line width=0.9pt,fill=FigureModel!9,minimum width=30mm,minimum height=8mm,align=center},
  bad tag/.style={tag,draw=BookAmber,fill=BookAmber!9},
  choices/.style={draw=BookMuted,line width=0.7pt,dashed,fill=white,minimum width=34mm,minimum height=9mm,align=center,font=\figurefont\scriptsize},
  allowed/.style={knowledge arrow,draw=BookInk,line width=1pt},
  invalid/.style={knowledge arrow,draw=BookAmber,line width=1pt,dashed},
  edge label/.style={fill=white,text=BookMuted,inner sep=1pt,font=\figurefont\scriptsize}
]
  \fill[white] (0,0) rectangle (16.85,6.3);
  \node[anchor=east,font=\figurefont\bfseries] at (1.25,5.85) {词元};
  \foreach \x/\word in {3.1/林舟,7.3/加入,11.5/海岬,15.1/科技} {
    \node[token] at (\x,5.85) {\word};
  }

  \node[anchor=east,font=\figurefont\bfseries,text=BookAmber] at (1.25,4.45) {局部};
  \node[tag] (bad1) at (3.1,4.45) {\texttt{B-PER}};
  \node[bad tag] (bad2) at (7.3,4.45) {\texttt{I-ORG}};
  \node[tag] (bad3) at (11.5,4.45) {\texttt{B-ORG}};
  \node[tag] (bad4) at (15.1,4.45) {\texttt{I-ORG}};
  \draw[invalid] (bad1) -- (bad2);
  \draw[allowed] (bad2) -- (bad3);
  \draw[allowed] (bad3) -- (bad4);
  \node[circle,draw=BookAmber,fill=white,text=BookAmber,line width=0.8pt,minimum size=4mm,inner sep=0pt,font=\figurefont\bfseries\scriptsize] at (5.2,4.45) {$\times$};

  \node[anchor=east,font=\figurefont\bfseries,text=FigureModel] at (1.25,3.0) {合法};
  \node[tag] (good1) at (3.1,3.0) {\texttt{B-PER}};
  \node[tag] (good2) at (7.3,3.0) {\texttt{O}};
  \node[tag] (good3) at (11.5,3.0) {\texttt{B-ORG}};
  \node[tag] (good4) at (15.1,3.0) {\texttt{I-ORG}};
  \draw[allowed] (good1) -- node[edge label,above] {允许} (good2);
  \draw[allowed] (good2) -- node[edge label,above] {允许} (good3);
  \draw[allowed] (good3) -- node[edge label,above] {允许} (good4);

  \node[anchor=east,font=\figurefont\bfseries,text=BookMuted] at (1.25,1.45) {候选};
  \node[choices] at (3.1,1.45) {\{\texttt{B-PER}, \texttt{O}\}};
  \node[choices] at (7.3,1.45) {\{\texttt{O}, \texttt{I-ORG}\}};
  \node[choices] at (11.5,1.45) {\{\texttt{B-ORG}, \texttt{O}\}};
  \node[choices] at (15.1,1.45) {\{\texttt{I-ORG}, \texttt{O}\}};
  \node[anchor=south,text=BookMuted,font=\figurefont\scriptsize,align=center] at (8.4,0.2) {结构约束排除非法组合，但合法序列仍可能不止一个};
\end{tikzpicture}

  \caption{局部标签与整体结构的约束关系。采用本节的IOB2边界约定：\texttt{B-X}开始类型X片段，\texttt{I-X}只能延续同类型的\texttt{B-X}或\texttt{I-X}，\texttt{O}位于片段外。上行的\texttt{B-PER}\,$\to$\,\texttt{I-ORG}违反转移条件；下行是一种合法候选。每个位置下方保留局部候选集合，联合约束筛选后仍可能留下多个合法序列。}
  \label{fig:rule-based-knowledge-acquisition-sequence-constraint}
\end{figure*}

结构约束不仅能事后验收，还能在求解过程中主动缩小候选，这就是\emph{局部传播}。其做法是：反复扫描相邻位置，凡是某个位置的候选标签在相邻位置中\emph{找不到任何合法搭配}，就把它删掉；不断重复，直到再也删不动为止。

在上面的序列里，第二位置的\texttt{I-ORG}在第一位置的候选集$\{\texttt{B-PER},\texttt{O}\}$中找不到任何一个能充当其合法前驱的标签（两者后面都不能直接跟\texttt{I-ORG}），因此它在第一轮就被删除。传播会持续进行：第四位置的\texttt{I-ORG}暂时不能删，因为第三位置还保留着\texttt{B-ORG}可以作它的前驱；但要注意这种支持是\emph{有条件}的——一旦第三位置最终选了\texttt{O}，第四位置的\texttt{I-ORG}就失去依托、不能再选。

这说明一个重要事实：传播结束后每个位置\emph{各自}仍留有候选，并不代表把这些候选\emph{任意}拼起来都是合法序列，跨位置的搭配约束依然存在。局部传播只是剪掉了“局部就注定不可能”的标签；要得到一个完整、全局合法的赋值，或要在所有合法序列中按语义评分挑出最优的一个，还必须进行\emph{联合选择}，根据结构形态选用动态规划（如序列）、搜索或通用约束求解器。

局部传播在链状序列上很有效，但把结构从序列换成一般的图时，“每处局部都可行”并不能保证“整体存在解”，一个经典例子是三节点的着色式矛盾。设三个变量$a,b,c$都只能取$0$或$1$，并要求\emph{每一对}变量取值都不同（即$a\ne b$、$b\ne c$、$c\ne a$）。

逐个局部检查会发现一切正常：对任意一对变量、任意一个取值，都能在另一个变量上找到相反值来满足“二者不同”，因此局部传播删不掉任何候选。但全局地看，两种赋值都会走进死胡同：若令$a=0$，为与$a$不同只能有$b=1$、$c=1$，可这样$b$与$c$相等，违反$b\ne c$；令$a=1$同理导致$b=c=0$，仍矛盾。结论是：三条两两约束\emph{局部都有支持}，合在一起却\emph{没有任何共同可行的赋值}（这正是奇环不能二染色的现象）。这提醒我们，局部一致性只是全局有解的必要条件而非充分条件，某些冲突只有在联合求解时才暴露。

结构约束的思想还可以扩展到更多关系：片段之间的\emph{边界包含}（大片段能否套小片段）、标签的\emph{父子层级}（子标签出现时父标签是否必须出现）、以及\emph{跨样本的身份一致性}（同一个实体在不同样本中应保持同一标识）。使用时都要把具体要检查的结构关系讲清楚，不能只说“加结构约束”。特别要注意身份一致与属性可变的区分：跨样本要求一致的是实体的\emph{标识}（这个人始终是同一个人），并\emph{不}要求他的岗位、位置或状态也一直不变——人会换工作、会移动，把身份一致错误地扩大成属性不变，会产生大量伪冲突。

所有这些结构推理都建立在一个前提上：那张“谁挨着谁、谁包含谁”的结构本身是有依据的。而结构可能来自可靠程度很不相同的渠道：文档层级可能直接来自文件格式（较可靠），跨样本的同一对象关系要靠实体匹配确定（可能出错），图中的邻接边则可能由距离阈值或某条关系规则构造出来（带有主观性）。

结构一旦有错，约束推理会“忠实地”沿着错误结构得出错误结论，且形式上毫无破绽：缺了一条本应存在的边，会阻断本该发生的传播，使两个本应联动的对象失去约束；多出一条错误的边，会把毫不相干的对象强行联结，逼它们满足本不该有的关系；边界定错，则可能把本来合法的标签组合误判成冲突。因此结构的来源与版本必须像标签证据一样保存；而且要清醒，约束求解\emph{成功}（得到了一个干净的合法解）绝不能反过来证明所用结构正确——结构错了，求解照样能成功。

结构约束的求解成本高度依赖结构的规模和连接方式，选择算法时要据此判断。对只含相邻转移条件的长序列，存在高效的动态规划解法（如Viterbi式递推），可以在与序列长度成线性或多项式的时间内求出全局最优合法序列；但当结构是连接稠密的图、约束又是全局性的（如三节点那种相互耦合），问题可能要求指数级的搜索空间，未必存在高效精确算法。

若因此改用近似方法，求解说明中必须如实写明两件事：还有哪些冲突没有真正解决，以及返回的结果到底是“满足全部约束的完整解”，还是“只满足了部分局部条件的近似结果”。把近似解包装成完整解，会让下游误以为所有结构约束都已满足。

\subsection{基于领域约束的规则利用}
\label{sec:rule-based-knowledge-acquisition-domain-constraint}

逻辑约束和结构约束具有较强的通用性，与具体行业无关；最后一类约束则直接来自某个\emph{具体领域的客观规律或业务规则}。一个直观例子是盘点物料：如果知道期初库存量、这段时间的流入量和期末库存量，那么那段没有直接测量的流出量其实被“物料守恒”限制住了——东西不会凭空产生或消失。

\term{领域约束}（domain constraint）把这种守恒关系、数量平衡、取值范围限制或业务一致性规则，写成未知标签或未知数值必须满足的条件，再用已有的观测去缩小未知量的候选范围。引入它并不意味着要换一个全新的求解器（很多情况下仍是在解同一类约束满足问题），真正改变的是\emph{约束的依据}：它来自特定领域的知识，因此必须额外讲清每个变量代表什么物理量、在什么环境与前提下成立、用什么单位、允许多大容差。脱离这些限定，守恒方程可能根本不适用。

先在理想条件下算一个能得到唯一值的例子。考虑同一个核算时段内的一组教学观测：初始库存$s_0=100\,\mathrm{kg}$，期间流入$q_{\mathrm{in}}=30\,\mathrm{kg}$，期末盘点得到结束量$s_1=108\,\mathrm{kg}$，未知的是期间流出量$q_{\mathrm{out}}$。

在一个\emph{封闭}前提下——除了已记录的流入、流出外，没有任何其他来源的进料，也没有损耗、失窃等未计量损失——物料在账面上应当严格平衡：期末量等于期初量加流入、再减流出。于是有精确的平衡关系
\begin{equation}
  s_1=s_0+q_{\mathrm{in}}-q_{\mathrm{out}},
  \label{eq:rule-based-knowledge-acquisition-balance}
\end{equation}
把三个已知量代入，$108=100+30-q_{\mathrm{out}}$，解得$q_{\mathrm{out}}=22\,\mathrm{kg}$。这里能得到唯一、确定的值，靠的是\emph{两个条件同时成立}：三个观测完整准确，以及“没有其他进出与损失”的封闭前提。这个前提一旦不成立（比如有未登记的损耗），就不能再沿用同一条等式作同样的推断，否则算出来的$22$只是“账面差额”，不等于真实流出。

现实中测量必有误差，要求等式严格成立往往不现实，此时问题应从“求唯一值”改成“判断哪些取值仍然与观测相容”。引入\emph{整体平衡残差}来量化偏离程度：它等于“观测到的期末量”与“按平衡关系算出的期末量”之差；对残差取绝对值，是为了同时限制“盘多了”和“盘少了”\emph{两个方向}的偏离，而不只是单侧。

假设各项观测的误差综合起来，使整体残差的绝对值不超过$2\,\mathrm{kg}$，那么不再要求等式精确成立，而只要求
\[
  \left|s_1-(s_0+q_{\mathrm{in}}-q_{\mathrm{out}})\right|
  \leq 2\,\mathrm{kg},
\]
仍代入同一组数值，即$|108-(100+30-q_{\mathrm{out}})|=|q_{\mathrm{out}}-22|\leq2$，规则只能推出一个可行区间$q_{\mathrm{out}}\in[20,24]\,\mathrm{kg}$。请注意：这时$22$虽然位于区间中央，但它\emph{不再是被确定的答案}，真实流出量可能是区间内任意值，不能把区间中点误报成精确结论。还要厘清$2\,\mathrm{kg}$这个数的口径：它约束的是\emph{合成后的整体残差}，并不是对三个观测各自误差界的分别声明，二者不能混用。

边界情形也要能识别：如果连流入量都缺失，已知量不足，原观测就不再能确定上述区间，应保留无法确定；若把期末量改成$s_1=140\,\mathrm{kg}$、同时仍要求流出非负$q_{\mathrm{out}}\geq0$，那么在“无其他来源”的前提下，账怎么算都平不了（期初加流入只有130，不可能有140），这是\emph{无解}，提示观测或前提冲突。因此求解器应能分别返回三种结果：唯一值、候选范围、观测冲突，并保存单位、容差与求解精度。

一条领域关系能不能当作“必须精确成立的硬等式”，取决于它的\emph{依据有多强}，可以区分三个层次。第一，若关系是按任务定义\emph{必然成立}的业务恒等式（例如复式记账里借贷必须相等、各互斥账项之和必须等于总量），它可以直接作为硬约束，违反即说明记账出错。第二，若是一条\emph{领域规律}（如物料守恒），但它依赖系统封闭、过程稳态或其他理想化条件，就必须把这些条件明确写进适用前提；前提不满足时，规律不能硬套。第三，若关系只是从有限样本中\emph{观察、拟合}出来的经验相关性，它通常更适合写成带误差的软约束，并保留估计误差和“条件变化需重新估计”的更新条款，而不能当成恒等式。

表~\ref{tab:rule-based-knowledge-acquisition-domain-constraint}固定同一个物料平衡例子，并排比较这三种依据各自允许推出多强的结论：精确等式给唯一值，带容差的理想规律给可行区间，经验关系只给偏好或异常候选。三者的数学表达式可能看起来很像，但证据强度和可用结论完全不同，绝不能因为“形式相似”就互换使用。

\begin{table*}[htbp]
  \centering
  \begingroup
  \setlength{\tabcolsep}{2.5pt}
  \setlength{\parindent}{0pt}
  \setlength{\RaggedRightParindent}{0pt}
  \renewcommand{\arraystretch}{1.12}
  \begin{tabularx}{\linewidth}{>{\raggedright\arraybackslash}p{23mm}>{\raggedright\arraybackslash}p{32mm}>{\raggedright\arraybackslash}p{25mm}>{\raggedright\arraybackslash}p{26mm}>{\raggedright\arraybackslash}p{32mm}>{\raggedright\arraybackslash}X}
    \toprule
    依据类型 & 前提 & 观测 & 误差或容差 & 可得结论 & 失效处理 \\
    \midrule
    按定义成立的关系 & 各量是完整且互斥的记账项 & $s_0=100$、$q_{\mathrm{in}}=30$、$s_1=108$，单位为$\mathrm{kg}$ & 无额外容差 & 精确值$q_{\mathrm{out}}=22\,\mathrm{kg}$ & 报告定义冲突或无解 \\
    理想条件下的领域规律 & 除已列流入、流出外，无其他物质通量或损失；观测误差已合成 & 同上；另核验核算边界及传感器版本 & 整体平衡残差绝对值不超过$2\,\mathrm{kg}$ & 可行范围$q_{\mathrm{out}}\in[20,24]\,\mathrm{kg}$ & 范围外停用；无法确定时弃权 \\
    样本归纳的经验关系 & 关系仅在指定设备、时间与样本范围内观察到 & 同类测量及规则拟合样本的版本 & 软代价与经核验的残差范围 & 偏好排序或异常候选 & 条件变化后重新估计 \\
    \bottomrule
  \end{tabularx}
  \endgroup
  \caption{领域约束的前提决定结论强度。表中数值均为单位一致的假设教学观测值；精确等式、带容差范围和经验软关系分别支持唯一值、可行区间和偏好候选，不能相互替代。}
  \label{tab:rule-based-knowledge-acquisition-domain-constraint}
\end{table*}

拿到一条候选的领域约束后，落到数据上还要先核对“公式里的变量”与“样本里的观测字段”是否真的表示同一个量——库存模型里的“流入”，是否对应数据表中那个字段的口径？字段名相近不代表物理含义相同。对由模型建议出来的关系尤其要警惕：模型很容易把“两个量高度相关”误写成“两者恒等”，而相关关系既无方向保证、也允许偏差，根本不具备守恒式的强度。

领域约束同样有时效性和环境依赖：当生产工艺、计量设备或业务流程改变后，原来“无额外物料通量、无损失”的封闭前提、误差范围甚至业务规则都可能不再成立，需要重新检查前提是否仍有效。因此开发和维护一套领域约束，工作内容远不止“把公式写出来”：还包括观测对齐、单位换算、误差建模、无解情形的解释以及版本管理，这些都是保证约束长期可用的必要成本。

最后把领域约束的几种结果按章首的状态约定逐一落成样本标注，并再次强调：平衡求解\emph{完成}这件事，本身并没有独立验证测量观测是否可靠——数学上解得出值，不代表传感器读数就是对的。

在“任务允许返回可行范围”的设定下：精确例提交唯一值$22\,\mathrm{kg}$，带容差例提交区间$[20,24]\,\mathrm{kg}$；两者的来源动作都是作答、执行成功、语义状态为已给出、验证状态为候选（尚未独立核验），区别在于后者给出的是范围、其中单个值仍未确定。再看两种“答不出”的不同情形：如果无法确认是否还存在其他物料通量（封闭前提是否成立都不知道），则检查虽完成但应弃权，内容为空、语义为未知，并记录“缺少核算边界证据”这一具体原因；如果进一步\emph{确认}存在未计入的通量，则是规则的适用前提不成立，该规则对此样本不适用，内容与动作都为空。最后是无解情形：当完整观测为$s_1=140\,\mathrm{kg}$、且非负流出要求不变时，求解成功地\emph{证明了无解}，但仍不提交流出量——动作是弃权、内容为空、语义为未知，同时把无解诊断和相互冲突的依据另行存档（这些诊断绝不能反过来被当成一个流出量标签）。

请特别区分最后这一对：“因前提不足而无法判断”与“已严格证明条件冲突而无解”，二者都导致没有任务赋值，但留下的原因和求解结果截然不同，必须分别记录；而“任务允许的多解范围”本身则是一项\emph{可以正式提交}的内容。规则产出的标签或约束质量究竟如何、错误怎样修正，将由第~\ref{chap:knowledge-quality-assessment-correction}章继续检查和处理。

\section{本章小结}
\label{sec:rule-based-knowledge-acquisition-summary}

回到章首的任命新闻可以看到，基于规则的知识获取绝不能止步于“程序在句子里找到了‘任命’这个词”。从一次字面命中到一个可用的候选标签，中间要走一条可追溯的链：先确认样本是否满足条件、再把提及关联到正确的对象、最后用答案之间的约束排除矛盾组合，而且每走一步都要说清依据是什么。

贯穿全章有两条彼此独立的主线，读者应当把它们分清。一条关于\emph{规则从哪来}：专家编写、从样本归纳、由模型生成，决定的是候选规则\emph{怎样被得到}。另一条关于\emph{规则靠什么下判断}：模式匹配、外部事实关联、约束推理，决定的是候选标签或约束\emph{凭什么产生}。这两条线可以自由交叉，一条规则也完全可以同时调用多种机制；但无论怎样组合，每一步都要分别保存待检查样本、规则结果、支持证据和版本信息，否则出错时无法沿链定位。

三类依据各有其能力边界，选用时应对各自的失效方式保持清醒。模式匹配不依赖外部资源、可以直接返回标签或关系候选，但它受符号变化、数值基准和条件组合的牵制，需要持续维护分词、阈值与分支逻辑。外部事实关联能够复用知识库里已有的实体、概念和关系，省去重复判断，但必须逐一检查实体消歧是否正确、事实版本与时间是否适用、以及证据范围是否真的覆盖当前句子——知识库为真不等于本句表达了它。约束推理不直接产生答案，而是排除不合法的组合、缩小候选范围或在无解时暴露冲突；它的一切结论都建立在前提和结构本身正确之上，结构错了，求解再干净也无济于事。尤其要记住，按软约束选出的“最小代价”候选，只代表它最符合\emph{当前人为设定的偏好}，这个选择动作本身并没有把它变成被证实的真实答案。

把任命例子中的各步串成一条完整链路，可以复习全章的方法：符号模式先在句中找到任命谓词及其参与者；实体链接接着把“林舟”“海岬科技”这两个提及对应到知识库中的唯一标识；关系匹配再按关系定义、方向、有效时间和证据范围，检查这个任职候选是否成立；当多个标签需要联合决定时，还可以用前面讲过的逻辑约束或结构约束筛掉矛盾组合。

这条链上所有的条件检查都遵循同一套三值约定，值得再集中复述一次：适用前提不成立时记为范围外、语义为不适用；在适用的前提下，前提与触发条件的合取为真才触发，为假则未触发且来源动作为空，无法确定时才弃权。合取遵循“假值优先于未知”——一个已确认的假值不应被另一个未知条件遮蔽。但强Kleene合取本身并不规定规则列表或规则树遇到未知时是否继续，那属于需要另行制定的\emph{控制策略}；触发之后若还有求解，求解结果也要单独记录。到了产出阶段：任务允许返回范围时，多个可行解可以作为候选内容正式作答；可行集为空说明无解，应如实报告冲突，而不是伪造一个标签。最后再强调全章反复出现的那条底线：\textbf{规则没有结果（范围外、未触发、弃权、未查到）都不等于负例}，一个真正的负标签仍然需要规则去明确检查负类的定义。

规则之所以值得使用，很大程度上是因为它的判断过程是\emph{显式}的：从待检查样本，到适用前提、模式匹配、事实关联，再到约束求解，每一步都留有记录，使错误可以沿链条被追查和定位——这是它相对“黑箱”判断的独特优势。但显式和可追溯\emph{并不}保证所有样本都能得到正确答案，规则会过时、会有覆盖不到的写法、也会依赖出错的前处理和外部事实。

因此是否继续使用一条规则，应当在三类因素之间权衡：它的依据是否仍可核验、它的适用范围是否还覆盖当前数据、以及维护它所需的前处理与更新成本是否划算。当缺少可用的规则条件或外部事实时，工作方向是补充覆盖；当多个来源（多条规则、规则与模型、规则与人）给出冲突结果时，不应在规则层强行裁决，而应进入第~\ref{chap:knowledge-fusion}章做多来源融合；而这些规则产出的标签和约束究竟可不可靠、发现错误后如何修正，则由第~\ref{chap:knowledge-quality-assessment-correction}章用统一的质量评估与修正流程收口。
